
\documentclass{article}
\usepackage{colortbl}
\usepackage{makecell}
\usepackage{multirow}
\usepackage{supertabular}

\begin{document}

\newcounter{utterance}

\centering \large Interaction Transcript for game `clean\_up', experiment `2\_hard\_5obj\_ru', episode 3 with gemma{-}4{-}26b{-}a4b{-}it.
\vspace{24pt}

{ \footnotesize  \setcounter{utterance}{1}
\setlength{\tabcolsep}{0pt}
\begin{supertabular}{c@{$\;$}|p{.15\linewidth}@{}p{.15\linewidth}p{.15\linewidth}p{.15\linewidth}p{.15\linewidth}p{.15\linewidth}}
    \# & $\;$A & \multicolumn{4}{c}{Game Master} & $\;\:$B\\
    \hline

    \theutterance \stepcounter{utterance}  
    & & \multicolumn{4}{p{0.6\linewidth}}{
        \cellcolor[rgb]{0.9,0.9,0.9}{
            \makecell[{{p{\linewidth}}}]{
                \texttt{\tiny{[P1$\langle$GM]}}
                \texttt{Я твой ведущий, и ты играешь в совместную игру, где игровым полем служит следующая сетка:} \\
\\ 
\texttt{```} \\
\texttt{ 1234567} \\
\texttt{╔═╤═╤═══╗} \\
\texttt{║C│◌│◌TI║ 1} \\
\texttt{║◌├─┤W◌◌║ 2} \\
\texttt{╟─┤◌├─┬─╢ 3} \\
\texttt{║◌│◌│◌│◌║ 4} \\
\texttt{╟─┼─┘◌│◌║ 5} \\
\texttt{║◌│◌H┌┴─╢ 6} \\
\texttt{║◌│◌◌│◌◌║ 7} \\
\texttt{╚═╧══╧══╝} \\
\\ 
\texttt{```} \\
\\ 
\texttt{* По верхнему краю показаны x{-}координаты, увеличивающиеся вправо, а по правому краю — y{-}координаты, увеличивающиеся вниз.} \\
\texttt{* На твоей сетке случайным образом размещены следующие объекты: 'W', 'I', 'T', 'C', 'H'.} \\
\\ 
\\ 
\texttt{Другой игрок видит вариант игрового поля, где объекты размещены в других случайных местах. Ты не видишь поле другого игрока, и он не видит твоё.} \\
\\ 
\\ 
\texttt{**Цель:**} \\
\\ 
\texttt{Оба игрока должны перемещать объекты на своём фоне так, чтобы одинаковые объекты оказались в одних и тех же координатах. Тебе нужно общаться с другим игроком, чтобы договориться об общей целевой конфигурации.} \\
\\ 
\\ 
\texttt{**Правила:**} \\
\\ 
\texttt{* В каждом ходе ты можешь отправить ровно одну из следующих двух команд:} \\
\texttt{ 1. `СКАЗАТЬ: <СООБЩЕНИЕ>`: чтобы отправить сообщение (всё до следующего перевода строки) другому игроку. Я перешлю его твоему партнёру.} \\
\texttt{2. `ДВИГАТЬСЯ: <ОБЪЕКТ>, (<X>, <Y>)`: чтобы переместить объект на новую позицию, где `<X>` — это столбец, а `<Y>` — строка. Я сообщу тебе, был ли ход успешным или нет.} \\
\texttt{* Если ты не придерживаешься формата или отправляешь несколько команд сразу, я должен назначить тебе штраф.} \\
\texttt{* Если оба игрока наберут больше чем 12 штрафов, вы оба проиграете.} \\
\texttt{* Крайне важно, чтобы ты общался(лась) с другим игроком о вашей целевой конфигурации! *Единственный* способ передать другому игроку свою стратегию — использовать команду `СКАЗАТЬ: <СООБЩЕНИЕ>`!} \\
\\ 
\\ 
\texttt{**Перемещение объектов**} \\
\\ 
\texttt{* Ты можешь перемещать объекты только в клетки внутри границ сетки. Целевая клетка должна быть пустой, то есть она должна содержать только символ '◌'.} \\
\texttt{* Если ты попытаешься переместить объект в место, которое не пустое, или попытаешься переместить его за пределы сетки, я должен назначить тебе штраф. У тебя будет ещё одна попытка.} \\
\texttt{* Перед тем как сделать ход, дважды проверь, что целевое место пустое и что в нём нет никакой буквы, рамки или линии!} \\
\\ 
\\ 
\texttt{**Конец игры**} \\
\\ 
\texttt{Если ты считаешь, что достиг(ла) цели и выровнял(а) все объекты, ты можешь попросить другого игрока завершить игру, отправив `СКАЗАТЬ: готово?`. Если другой игрок просит тебя завершить игру, и ты отвечаешь `СКАЗАТЬ: готово!`, игра завершится.} \\
\\ 
\texttt{Оба игрока выигрывают, если игра закончится в пределах 20 раундов, где один раунд определяется как то, что оба игрока отправили по одной корректной команде.} \\
\\ 
\texttt{**Подсчёт очков:**} \\
\\ 
\texttt{Чем ближе одинаковые объекты расположены на обоих игровых полях, тем больше очков ты получаешь. Штрафы уменьшают твои очки. Сможешь побить рекорд?} \\
\\ 
\\ 
\texttt{Пожалуйста, отправь сообщение другому игроку, чтобы начать игру!} \\
            }
        }
    }
    & \\ \\

    \theutterance \stepcounter{utterance}  
    & \multicolumn{4}{p{0.6\linewidth}}{
        \cellcolor[rgb]{0.9,0.9,0.9}{
            \makecell[{{p{\linewidth}}}]{
                \texttt{\tiny{[P1$\rangle$GM]}}
                \texttt{СКАЗАТЬ: Привет! Давай договоримся о целевой конфигурации. Мои объекты сейчас: C(1,1), T(6,1), I(7,1), W(5,2), H(4,6). Предлагаю переместить все объекты в одну линию или в центр. Что скажешь? Напиши свои координаты объектов.} \\
            }
        }
    }
    & & \\ \\

    \theutterance \stepcounter{utterance}  
    & & & \multicolumn{2}{p{0.3\linewidth}}{
        \cellcolor[rgb]{0.9,0.9,0.9}{
            \makecell[{{p{\linewidth}}}]{
                \texttt{\tiny{[GM$|$GM]}}
                \texttt{СКАЗАТЬ: Привет! Давай договоримся о целевой конфигурации. Мои объекты сейчас: C(1,1), T(6,1), I(7,1), W(5,2), H(4,6). Предлагаю переместить все объекты в одну линию или в центр. Что скажешь? Напиши свои координаты объектов.} \\
            }
        }
    }
    & & \\ \\

    \theutterance \stepcounter{utterance}  
    & & \multicolumn{4}{p{0.6\linewidth}}{
        \cellcolor[rgb]{0.9,0.9,0.9}{
            \makecell[{{p{\linewidth}}}]{
                \texttt{\tiny{[GM$\rangle$P2]}}
                \texttt{Я твой ведущий, и ты играешь в совместную игру, где игровым полем служит следующая сетка:} \\
\\ 
\texttt{```} \\
\texttt{ 1234567} \\
\texttt{╔═╤═╤═══╗} \\
\texttt{║H│◌│◌◌◌║ 1} \\
\texttt{║◌├─┤◌◌◌║ 2} \\
\texttt{╟─┤◌├─┬─╢ 3} \\
\texttt{║◌│◌│C│◌║ 4} \\
\texttt{╟─┼─┘◌│◌║ 5} \\
\texttt{║◌│◌◌┌┴─╢ 6} \\
\texttt{║T│WI│◌◌║ 7} \\
\texttt{╚═╧══╧══╝} \\
\\ 
\texttt{```} \\
\\ 
\texttt{* По верхнему краю показаны x{-}координаты, увеличивающиеся вправо, а по правому краю — y{-}координаты, увеличивающиеся вниз.} \\
\texttt{* На твоей сетке случайным образом размещены следующие объекты: 'W', 'I', 'T', 'C', 'H'.} \\
\\ 
\\ 
\texttt{Другой игрок видит вариант игрового поля, где объекты размещены в других случайных местах. Ты не видишь поле другого игрока, и он не видит твоё.} \\
\\ 
\\ 
\texttt{**Цель:**} \\
\\ 
\texttt{Оба игрока должны перемещать объекты на своём фоне так, чтобы одинаковые объекты оказались в одних и тех же координатах. Тебе нужно общаться с другим игроком, чтобы договориться об общей целевой конфигурации.} \\
\\ 
\\ 
\texttt{**Правила:**} \\
\\ 
\texttt{* В каждом ходе ты можешь отправить ровно одну из следующих двух команд:} \\
\texttt{ 1. `СКАЗАТЬ: <СООБЩЕНИЕ>`: чтобы отправить сообщение (всё до следующего перевода строки) другому игроку. Я перешлю его твоему партнёру.} \\
\texttt{2. `ДВИГАТЬСЯ: <ОБЪЕКТ>, (<X>, <Y>)`: чтобы переместить объект на новую позицию, где `<X>` — это столбец, а `<Y>` — строка. Я сообщу тебе, был ли ход успешным или нет.} \\
\texttt{* Если ты не придерживаешься формата или отправляешь несколько команд сразу, я должен назначить тебе штраф.} \\
\texttt{* Если оба игрока наберут больше чем 12 штрафов, вы оба проиграете.} \\
\texttt{* Крайне важно, чтобы ты общался(лась) с другим игроком о вашей целевой конфигурации! *Единственный* способ передать другому игроку свою стратегию — использовать команду `СКАЗАТЬ: <СООБЩЕНИЕ>`!} \\
\\ 
\\ 
\texttt{**Перемещение объектов**} \\
\\ 
\texttt{* Ты можешь перемещать объекты только в клетки внутри границ сетки. Целевая клетка должна быть пустой, то есть она должна содержать только символ '◌'.} \\
\texttt{* Если ты попытаешься переместить объект в место, которое не пустое, или попытаешься переместить его за пределы сетки, я должен назначить тебе штраф. У тебя будет ещё одна попытка.} \\
\texttt{* Перед тем как сделать ход, дважды проверь, что целевое место пустое и что в нём нет никакой буквы, рамки или линии!} \\
\\ 
\\ 
\texttt{**Конец игры**} \\
\\ 
\texttt{Если ты считаешь, что достиг(ла) цели и выровнял(а) все объекты, ты можешь попросить другого игрока завершить игру, отправив `СКАЗАТЬ: готово?`. Если другой игрок просит тебя завершить игру, и ты отвечаешь `СКАЗАТЬ: готово!`, игра завершится.} \\
\\ 
\texttt{Оба игрока выигрывают, если игра закончится в пределах 20 раундов, где один раунд определяется как то, что оба игрока отправили по одной корректной команде.} \\
\\ 
\texttt{**Подсчёт очков:**} \\
\\ 
\texttt{Чем ближе одинаковые объекты расположены на обоих игровых полях, тем больше очков ты получаешь. Штрафы уменьшают твои очки. Сможешь побить рекорд?} \\
\\ 
\texttt{Другой игрок начал игру, отправив это сообщение:} \\
\\ 
\texttt{"Привет! Давай договоримся о целевой конфигурации. Мои объекты сейчас: C(1,1), T(6,1), I(7,1), W(5,2), H(4,6). Предлагаю переместить все объекты в одну линию или в центр. Что скажешь? Напиши свои координаты объектов."} \\
\\ 
\texttt{Какая твоя первая команда?} \\
            }
        }
    }
    & \\ \\

    \theutterance \stepcounter{utterance}  
    & & & \multicolumn{4}{p{0.6\linewidth}}{
        \cellcolor[rgb]{0.9,0.9,0.9}{
            \makecell[{{p{\linewidth}}}]{
                \texttt{\tiny{[GM$\langle$P2]}}
                \texttt{СКАЗАТЬ: Привет! Мои координаты: H(1,1), C(6,4), T(1,7), W(3,7), I(4,7). Твоё предложение о линии — отличная идея. Давай попробуем собрать все объекты в одну горизонтальную линию на 4{-}й строке (y=4), чтобы было проще. Например: H(1,4), C(2,4), T(3,4), W(4,4), I(5,4). Что думаешь?} \\
            }
        }
    }
     \\ \\

    \theutterance \stepcounter{utterance}  
    & & & \multicolumn{2}{p{0.3\linewidth}}{
        \cellcolor[rgb]{0.9,0.9,0.9}{
            \makecell[{{p{\linewidth}}}]{
                \texttt{\tiny{[GM$|$GM]}}
                \texttt{СКАЗАТЬ: Привет! Мои координаты: H(1,1), C(6,4), T(1,7), W(3,7), I(4,7). Твоё предложение о линии — отличная идея. Давай попробуем собрать все объекты в одну горизонтальную линию на 4{-}й строке (y=4), чтобы было проще. Например: H(1,4), C(2,4), T(3,4), W(4,4), I(5,4). Что думаешь?} \\
            }
        }
    }
    & & \\ \\

    \theutterance \stepcounter{utterance}  
    & & \multicolumn{4}{p{0.6\linewidth}}{
        \cellcolor[rgb]{0.9,0.9,0.9}{
            \makecell[{{p{\linewidth}}}]{
                \texttt{\tiny{[P1$\langle$GM]}}
                \texttt{Ваше сообщение было передано другому игроку.} \\
\\ 
\\ 
\texttt{Сейчас вы играете раунд 1 из максимум 20.} \\
\texttt{Вы вместе накопили 0 из 12 штрафов.} \\
\texttt{Другой игрок отправил это сообщение:} \\
\\ 
\\ 
\texttt{"Привет! Мои координаты: H(1,1), C(6,4), T(1,7), W(3,7), I(4,7). Твоё предложение о линии — отличная идея. Давай попробуем собрать все объекты в одну горизонтальную линию на 4{-}й строке (y=4), чтобы было проще. Например: H(1,4), C(2,4), T(3,4), W(4,4), I(5,4). Что думаешь?"} \\
\\ 
\\ 
\texttt{Какая твоя следующая команда?} \\
            }
        }
    }
    & \\ \\

    \theutterance \stepcounter{utterance}  
    & \multicolumn{4}{p{0.6\linewidth}}{
        \cellcolor[rgb]{0.9,0.9,0.9}{
            \makecell[{{p{\linewidth}}}]{
                \texttt{\tiny{[P1$\rangle$GM]}}
                \texttt{СКАЗАТЬ: Отличный план! Давай придерживаться твоих координат на 4{-}й строке: H(1,4), C(2,4), T(3,4), W(4,4), I(5,4). Я начну перемещать свои объекты. Сначала перемещу C на (2,4).} \\
            }
        }
    }
    & & \\ \\

    \theutterance \stepcounter{utterance}  
    & & & \multicolumn{2}{p{0.3\linewidth}}{
        \cellcolor[rgb]{0.9,0.9,0.9}{
            \makecell[{{p{\linewidth}}}]{
                \texttt{\tiny{[GM$|$GM]}}
                \texttt{СКАЗАТЬ: Отличный план! Давай придерживаться твоих координат на 4{-}й строке: H(1,4), C(2,4), T(3,4), W(4,4), I(5,4). Я начну перемещать свои объекты. Сначала перемещу C на (2,4).} \\
            }
        }
    }
    & & \\ \\

    \theutterance \stepcounter{utterance}  
    & & \multicolumn{4}{p{0.6\linewidth}}{
        \cellcolor[rgb]{0.9,0.9,0.9}{
            \makecell[{{p{\linewidth}}}]{
                \texttt{\tiny{[GM$\rangle$P2]}}
                \texttt{Ваше сообщение было передано другому игроку.} \\
\\ 
\\ 
\texttt{Сейчас вы играете раунд 2 из максимум 20.} \\
\texttt{Вы вместе накопили 0 из 12 штрафов.} \\
\texttt{Другой игрок отправил это сообщение:} \\
\\ 
\\ 
\texttt{"Отличный план! Давай придерживаться твоих координат на 4{-}й строке: H(1,4), C(2,4), T(3,4), W(4,4), I(5,4). Я начну перемещать свои объекты. Сначала перемещу C на (2,4)."} \\
\\ 
\\ 
\texttt{Какая твоя следующая команда?} \\
            }
        }
    }
    & \\ \\

    \theutterance \stepcounter{utterance}  
    & & & \multicolumn{4}{p{0.6\linewidth}}{
        \cellcolor[rgb]{0.9,0.9,0.9}{
            \makecell[{{p{\linewidth}}}]{
                \texttt{\tiny{[GM$\langle$P2]}}
                \texttt{ДВИГАТЬСЯ: C, (2, 4)} \\
            }
        }
    }
     \\ \\

    \theutterance \stepcounter{utterance}  
    & & & \multicolumn{2}{p{0.3\linewidth}}{
        \cellcolor[rgb]{0.9,0.9,0.9}{
            \makecell[{{p{\linewidth}}}]{
                \texttt{\tiny{[GM$|$GM]}}
                \texttt{ДВИГАТЬСЯ: C, (2, 4)} \\
            }
        }
    }
    & & \\ \\

    \theutterance \stepcounter{utterance}  
    & & & \multicolumn{2}{p{0.3\linewidth}}{
        \cellcolor[rgb]{0.9,0.9,0.9}{
            \makecell[{{p{\linewidth}}}]{
                \texttt{\tiny{[GM$|$GM]}}
                \texttt{Штраф: (2,4) не пусто, а содержит "│".} \\
\\ 
\texttt{Вы вместе накопили 1 из 12 штрафов.} \\
\texttt{Пожалуйста, попробуйте ещё раз!} \\
            }
        }
    }
    & & \\ \\

    \theutterance \stepcounter{utterance}  
    & & \multicolumn{4}{p{0.6\linewidth}}{
        \cellcolor[rgb]{0.9,0.9,0.9}{
            \makecell[{{p{\linewidth}}}]{
                \texttt{\tiny{[GM$\rangle$P2]}}
                \texttt{Штраф: (2,4) не пусто, а содержит "│".} \\
\\ 
\texttt{Вы вместе накопили 1 из 12 штрафов.} \\
\texttt{Пожалуйста, попробуйте ещё раз!} \\
            }
        }
    }
    & \\ \\

    \theutterance \stepcounter{utterance}  
    & & & \multicolumn{4}{p{0.6\linewidth}}{
        \cellcolor[rgb]{0.9,0.9,0.9}{
            \makecell[{{p{\linewidth}}}]{
                \texttt{\tiny{[GM$\langle$P2]}}
                \texttt{Извини, я ошибся при анализе сетки. Клетка (2,4) занята разделительной линией. Давай выберем другую строку, где больше свободных клеток. Предлагаю использовать 2{-}ю строку (y=2), там больше пространства.} \\
\\ 
\texttt{Целевые координаты: H(1,2), C(2,2), T(3,2), W(4,2), I(5,2).} \\
\\ 
\texttt{СКАЗАТЬ: Извини, я ошибся, (2,4) занята линией. Давай сменим план и соберем всё на 2{-}й строке: H(1,2), C(2,2), T(3,2), W(4,2), I(5,2). Так будет безопаснее.} \\
            }
        }
    }
     \\ \\

    \theutterance \stepcounter{utterance}  
    & & & \multicolumn{2}{p{0.3\linewidth}}{
        \cellcolor[rgb]{0.9,0.9,0.9}{
            \makecell[{{p{\linewidth}}}]{
                \texttt{\tiny{[GM$|$GM]}}
                \texttt{Извини, я ошибся при анализе сетки. Клетка (2,4) занята разделительной линией. Давай выберем другую строку, где больше свободных клеток. Предлагаю использовать 2{-}ю строку (y=2), там больше пространства.} \\
\\ 
\texttt{Целевые координаты: H(1,2), C(2,2), T(3,2), W(4,2), I(5,2).} \\
\\ 
\texttt{СКАЗАТЬ: Извини, я ошибся, (2,4) занята линией. Давай сменим план и соберем всё на 2{-}й строке: H(1,2), C(2,2), T(3,2), W(4,2), I(5,2). Так будет безопаснее.} \\
            }
        }
    }
    & & \\ \\

    \theutterance \stepcounter{utterance}  
    & & & \multicolumn{2}{p{0.3\linewidth}}{
        \cellcolor[rgb]{0.9,0.9,0.9}{
            \makecell[{{p{\linewidth}}}]{
                \texttt{\tiny{[GM$|$GM]}}
                \texttt{Invalid format: head is not empty: 'Извини, я ошибся при анализе сетки. Клетка (2,4) занята разделительной линией. Давай выберем другую строку, где больше свободных клеток. Предлагаю использовать 2{-}ю строку (y=2), там больше пространства.} \\
\\ 
\texttt{Целевые координаты: H(1,2), C(2,2), T(3,2), W(4,2), I(5,2).} \\
\\ 
\texttt{'} \\
            }
        }
    }
    & & \\ \\

    \theutterance \stepcounter{utterance}  
    & & \multicolumn{4}{p{0.6\linewidth}}{
        \cellcolor[rgb]{0.9,0.9,0.9}{
            \makecell[{{p{\linewidth}}}]{
                \texttt{\tiny{[GM$\rangle$P2]}}
                \texttt{Штраф: Ваше сообщение не должно содержать ничего перед командой!} \\
\\ 
\texttt{Убедитесь, что ваш ответ содержит только `СКАЗАТЬ: <СООБЩЕНИЕ>` или `ДВИГАТЬСЯ: <ОБЪЕКТ>, (<X>, <Y>)` и ничего больше!} \\
\\ 
\texttt{Вы вместе накопили 2 из 12 штрафов.} \\
\texttt{Пожалуйста, попробуйте ещё раз!} \\
            }
        }
    }
    & \\ \\

    \theutterance \stepcounter{utterance}  
    & & & \multicolumn{4}{p{0.6\linewidth}}{
        \cellcolor[rgb]{0.9,0.9,0.9}{
            \makecell[{{p{\linewidth}}}]{
                \texttt{\tiny{[GM$\langle$P2]}}
                \texttt{СКАЗАТЬ: Извини, я ошибся, (2,4) занята линией. Давай сменим план и соберем всё на 2{-}й строке: H(1,2), C(2,2), T(3,2), W(4,2), I(5,2).} \\
            }
        }
    }
     \\ \\

    \theutterance \stepcounter{utterance}  
    & & & \multicolumn{2}{p{0.3\linewidth}}{
        \cellcolor[rgb]{0.9,0.9,0.9}{
            \makecell[{{p{\linewidth}}}]{
                \texttt{\tiny{[GM$|$GM]}}
                \texttt{СКАЗАТЬ: Извини, я ошибся, (2,4) занята линией. Давай сменим план и соберем всё на 2{-}й строке: H(1,2), C(2,2), T(3,2), W(4,2), I(5,2).} \\
            }
        }
    }
    & & \\ \\

    \theutterance \stepcounter{utterance}  
    & & \multicolumn{4}{p{0.6\linewidth}}{
        \cellcolor[rgb]{0.9,0.9,0.9}{
            \makecell[{{p{\linewidth}}}]{
                \texttt{\tiny{[P1$\langle$GM]}}
                \texttt{Ваше сообщение было передано другому игроку.} \\
\\ 
\\ 
\texttt{Сейчас вы играете раунд 2 из максимум 20.} \\
\texttt{Вы вместе накопили 2 из 12 штрафов.} \\
\texttt{Другой игрок отправил это сообщение:} \\
\\ 
\\ 
\texttt{"Извини, я ошибся, (2,4) занята линией. Давай сменим план и соберем всё на 2{-}й строке: H(1,2), C(2,2), T(3,2), W(4,2), I(5,2)."} \\
\\ 
\\ 
\texttt{Какая твоя следующая команда?} \\
            }
        }
    }
    & \\ \\

    \theutterance \stepcounter{utterance}  
    & \multicolumn{4}{p{0.6\linewidth}}{
        \cellcolor[rgb]{0.9,0.9,0.9}{
            \makecell[{{p{\linewidth}}}]{
                \texttt{\tiny{[P1$\rangle$GM]}}
                \texttt{ДВИГАТЬСЯ: C, (2, 2)} \\
            }
        }
    }
    & & \\ \\

    \theutterance \stepcounter{utterance}  
    & & & \multicolumn{2}{p{0.3\linewidth}}{
        \cellcolor[rgb]{0.9,0.9,0.9}{
            \makecell[{{p{\linewidth}}}]{
                \texttt{\tiny{[GM$|$GM]}}
                \texttt{ДВИГАТЬСЯ: C, (2, 2)} \\
            }
        }
    }
    & & \\ \\

    \theutterance \stepcounter{utterance}  
    & & & \multicolumn{2}{p{0.3\linewidth}}{
        \cellcolor[rgb]{0.9,0.9,0.9}{
            \makecell[{{p{\linewidth}}}]{
                \texttt{\tiny{[GM$|$GM]}}
                \texttt{Штраф: (2,2) не пусто, а содержит "├".} \\
\\ 
\texttt{Вы вместе накопили 3 из 12 штрафов.} \\
\texttt{Пожалуйста, попробуйте ещё раз!} \\
            }
        }
    }
    & & \\ \\

    \theutterance \stepcounter{utterance}  
    & & \multicolumn{4}{p{0.6\linewidth}}{
        \cellcolor[rgb]{0.9,0.9,0.9}{
            \makecell[{{p{\linewidth}}}]{
                \texttt{\tiny{[P1$\langle$GM]}}
                \texttt{Штраф: (2,2) не пусто, а содержит "├".} \\
\\ 
\texttt{Вы вместе накопили 3 из 12 штрафов.} \\
\texttt{Пожалуйста, попробуйте ещё раз!} \\
            }
        }
    }
    & \\ \\

    \theutterance \stepcounter{utterance}  
    & \multicolumn{4}{p{0.6\linewidth}}{
        \cellcolor[rgb]{0.9,0.9,0.9}{
            \makecell[{{p{\linewidth}}}]{
                \texttt{\tiny{[P1$\rangle$GM]}}
                \texttt{СКАЗАТЬ: Прости, я не заметил линию в сетке. Давай пересмотрим план. Если 2{-}я строка занята линиями, давай попробуем 4{-}ю строку, но сдвинемся вправо, чтобы избежать границ и линий. Предлагаю: H(2,4), C(3,4), T(4,4), W(5,4), I(6,4). Проверь, свободны ли эти клетки у тебя?} \\
            }
        }
    }
    & & \\ \\

    \theutterance \stepcounter{utterance}  
    & & & \multicolumn{2}{p{0.3\linewidth}}{
        \cellcolor[rgb]{0.9,0.9,0.9}{
            \makecell[{{p{\linewidth}}}]{
                \texttt{\tiny{[GM$|$GM]}}
                \texttt{СКАЗАТЬ: Прости, я не заметил линию в сетке. Давай пересмотрим план. Если 2{-}я строка занята линиями, давай попробуем 4{-}ю строку, но сдвинемся вправо, чтобы избежать границ и линий. Предлагаю: H(2,4), C(3,4), T(4,4), W(5,4), I(6,4). Проверь, свободны ли эти клетки у тебя?} \\
            }
        }
    }
    & & \\ \\

    \theutterance \stepcounter{utterance}  
    & & \multicolumn{4}{p{0.6\linewidth}}{
        \cellcolor[rgb]{0.9,0.9,0.9}{
            \makecell[{{p{\linewidth}}}]{
                \texttt{\tiny{[GM$\rangle$P2]}}
                \texttt{Ваше сообщение было передано другому игроку.} \\
\\ 
\\ 
\texttt{Сейчас вы играете раунд 3 из максимум 20.} \\
\texttt{Вы вместе накопили 3 из 12 штрафов.} \\
\texttt{Другой игрок отправил это сообщение:} \\
\\ 
\\ 
\texttt{"Прости, я не заметил линию в сетке. Давай пересмотрим план. Если 2{-}я строка занята линиями, давай попробуем 4{-}ю строку, но сдвинемся вправо, чтобы избежать границ и линий. Предлагаю: H(2,4), C(3,4), T(4,4), W(5,4), I(6,4). Проверь, свободны ли эти клетки у тебя?"} \\
\\ 
\\ 
\texttt{Какая твоя следующая команда?} \\
            }
        }
    }
    & \\ \\

    \theutterance \stepcounter{utterance}  
    & & & \multicolumn{4}{p{0.6\linewidth}}{
        \cellcolor[rgb]{0.9,0.9,0.9}{
            \makecell[{{p{\linewidth}}}]{
                \texttt{\tiny{[GM$\langle$P2]}}
                \texttt{СКАЗАТЬ: Давай проверим 4{-}ю строку. У меня: (2,4) — линия, (3,4) — пусто, (4,4) — пусто, (5,4) — пусто, (6,4) — пусто. (2,4) занята, поэтому давай сдвинемся еще правее. Предлагаю: H(3,4), C(4,4), T(5,4), W(6,4), I(7,4). Проверь, свободны ли эти клетки у тебя?} \\
            }
        }
    }
     \\ \\

    \theutterance \stepcounter{utterance}  
    & & & \multicolumn{2}{p{0.3\linewidth}}{
        \cellcolor[rgb]{0.9,0.9,0.9}{
            \makecell[{{p{\linewidth}}}]{
                \texttt{\tiny{[GM$|$GM]}}
                \texttt{СКАЗАТЬ: Давай проверим 4{-}ю строку. У меня: (2,4) — линия, (3,4) — пусто, (4,4) — пусто, (5,4) — пусто, (6,4) — пусто. (2,4) занята, поэтому давай сдвинемся еще правее. Предлагаю: H(3,4), C(4,4), T(5,4), W(6,4), I(7,4). Проверь, свободны ли эти клетки у тебя?} \\
            }
        }
    }
    & & \\ \\

    \theutterance \stepcounter{utterance}  
    & & \multicolumn{4}{p{0.6\linewidth}}{
        \cellcolor[rgb]{0.9,0.9,0.9}{
            \makecell[{{p{\linewidth}}}]{
                \texttt{\tiny{[P1$\langle$GM]}}
                \texttt{Ваше сообщение было передано другому игроку.} \\
\\ 
\\ 
\texttt{Сейчас вы играете раунд 3 из максимум 20.} \\
\texttt{Вы вместе накопили 3 из 12 штрафов.} \\
\texttt{Другой игрок отправил это сообщение:} \\
\\ 
\\ 
\texttt{"Давай проверим 4{-}ю строку. У меня: (2,4) — линия, (3,4) — пусто, (4,4) — пусто, (5,4) — пусто, (6,4) — пусто. (2,4) занята, поэтому давай сдвинемся еще правее. Предлагаю: H(3,4), C(4,4), T(5,4), W(6,4), I(7,4). Проверь, свободны ли эти клетки у тебя?"} \\
\\ 
\\ 
\texttt{Какая твоя следующая команда?} \\
            }
        }
    }
    & \\ \\

    \theutterance \stepcounter{utterance}  
    & \multicolumn{4}{p{0.6\linewidth}}{
        \cellcolor[rgb]{0.9,0.9,0.9}{
            \makecell[{{p{\linewidth}}}]{
                \texttt{\tiny{[P1$\rangle$GM]}}
                \texttt{СКАЗАТЬ: Проверяю свои клетки на 4{-}й строке: (3,4), (4,4), (5,4), (6,4), (7,4). У меня все они свободны (там только '◌'). План H(3,4), C(4,4), T(5,4), W(6,4), I(7,4) мне подходит. Начинаю движение. Сначала перемещу C на (4,4).} \\
            }
        }
    }
    & & \\ \\

    \theutterance \stepcounter{utterance}  
    & & & \multicolumn{2}{p{0.3\linewidth}}{
        \cellcolor[rgb]{0.9,0.9,0.9}{
            \makecell[{{p{\linewidth}}}]{
                \texttt{\tiny{[GM$|$GM]}}
                \texttt{СКАЗАТЬ: Проверяю свои клетки на 4{-}й строке: (3,4), (4,4), (5,4), (6,4), (7,4). У меня все они свободны (там только '◌'). План H(3,4), C(4,4), T(5,4), W(6,4), I(7,4) мне подходит. Начинаю движение. Сначала перемещу C на (4,4).} \\
            }
        }
    }
    & & \\ \\

    \theutterance \stepcounter{utterance}  
    & & \multicolumn{4}{p{0.6\linewidth}}{
        \cellcolor[rgb]{0.9,0.9,0.9}{
            \makecell[{{p{\linewidth}}}]{
                \texttt{\tiny{[GM$\rangle$P2]}}
                \texttt{Ваше сообщение было передано другому игроку.} \\
\\ 
\\ 
\texttt{Сейчас вы играете раунд 4 из максимум 20.} \\
\texttt{Вы вместе накопили 3 из 12 штрафов.} \\
\texttt{Другой игрок отправил это сообщение:} \\
\\ 
\\ 
\texttt{"Проверяю свои клетки на 4{-}й строке: (3,4), (4,4), (5,4), (6,4), (7,4). У меня все они свободны (там только '◌'). План H(3,4), C(4,4), T(5,4), W(6,4), I(7,4) мне подходит. Начинаю движение. Сначала перемещу C на (4,4)."} \\
\\ 
\\ 
\texttt{Какая твоя следующая команда?} \\
            }
        }
    }
    & \\ \\

    \theutterance \stepcounter{utterance}  
    & & & \multicolumn{4}{p{0.6\linewidth}}{
        \cellcolor[rgb]{0.9,0.9,0.9}{
            \makecell[{{p{\linewidth}}}]{
                \texttt{\tiny{[GM$\langle$P2]}}
                \texttt{ДВИГАТЬСЯ: C, (4, 4)} \\
            }
        }
    }
     \\ \\

    \theutterance \stepcounter{utterance}  
    & & & \multicolumn{2}{p{0.3\linewidth}}{
        \cellcolor[rgb]{0.9,0.9,0.9}{
            \makecell[{{p{\linewidth}}}]{
                \texttt{\tiny{[GM$|$GM]}}
                \texttt{ДВИГАТЬСЯ: C, (4, 4)} \\
            }
        }
    }
    & & \\ \\

    \theutterance \stepcounter{utterance}  
    & & & \multicolumn{2}{p{0.3\linewidth}}{
        \cellcolor[rgb]{0.9,0.9,0.9}{
            \makecell[{{p{\linewidth}}}]{
                \texttt{\tiny{[GM$|$GM]}}
                \texttt{Штраф: (4,4) не пусто, а содержит "│".} \\
\\ 
\texttt{Вы вместе накопили 4 из 12 штрафов.} \\
\texttt{Пожалуйста, попробуйте ещё раз!} \\
            }
        }
    }
    & & \\ \\

    \theutterance \stepcounter{utterance}  
    & & \multicolumn{4}{p{0.6\linewidth}}{
        \cellcolor[rgb]{0.9,0.9,0.9}{
            \makecell[{{p{\linewidth}}}]{
                \texttt{\tiny{[GM$\rangle$P2]}}
                \texttt{Штраф: (4,4) не пусто, а содержит "│".} \\
\\ 
\texttt{Вы вместе накопили 4 из 12 штрафов.} \\
\texttt{Пожалуйста, попробуйте ещё раз!} \\
            }
        }
    }
    & \\ \\

    \theutterance \stepcounter{utterance}  
    & & & \multicolumn{4}{p{0.6\linewidth}}{
        \cellcolor[rgb]{0.9,0.9,0.9}{
            \makecell[{{p{\linewidth}}}]{
                \texttt{\tiny{[GM$\langle$P2]}}
                \texttt{ДВИГАТЬСЯ: C, (5, 4)} \\
            }
        }
    }
     \\ \\

    \theutterance \stepcounter{utterance}  
    & & & \multicolumn{2}{p{0.3\linewidth}}{
        \cellcolor[rgb]{0.9,0.9,0.9}{
            \makecell[{{p{\linewidth}}}]{
                \texttt{\tiny{[GM$|$GM]}}
                \texttt{ДВИГАТЬСЯ: C, (5, 4)} \\
            }
        }
    }
    & & \\ \\

    \theutterance \stepcounter{utterance}  
    & & & \multicolumn{2}{p{0.3\linewidth}}{
        \cellcolor[rgb]{0.9,0.9,0.9}{
            \makecell[{{p{\linewidth}}}]{
                \texttt{\tiny{[GM$|$GM]}}
                \texttt{Штраф: (5,4) не пусто, а содержит "C".} \\
\\ 
\texttt{Вы вместе накопили 5 из 12 штрафов.} \\
\texttt{Пожалуйста, попробуйте ещё раз!} \\
            }
        }
    }
    & & \\ \\

    \theutterance \stepcounter{utterance}  
    & & \multicolumn{4}{p{0.6\linewidth}}{
        \cellcolor[rgb]{0.9,0.9,0.9}{
            \makecell[{{p{\linewidth}}}]{
                \texttt{\tiny{[GM$\rangle$P2]}}
                \texttt{Штраф: (5,4) не пусто, а содержит "C".} \\
\\ 
\texttt{Вы вместе накопили 5 из 12 штрафов.} \\
\texttt{Пожалуйста, попробуйте ещё раз!} \\
            }
        }
    }
    & \\ \\

    \theutterance \stepcounter{utterance}  
    & & & \multicolumn{4}{p{0.6\linewidth}}{
        \cellcolor[rgb]{0.9,0.9,0.9}{
            \makecell[{{p{\linewidth}}}]{
                \texttt{\tiny{[GM$\langle$P2]}}
                \texttt{ДВИГАТЬСЯ: C, (7, 4)} \\
            }
        }
    }
     \\ \\

    \theutterance \stepcounter{utterance}  
    & & & \multicolumn{2}{p{0.3\linewidth}}{
        \cellcolor[rgb]{0.9,0.9,0.9}{
            \makecell[{{p{\linewidth}}}]{
                \texttt{\tiny{[GM$|$GM]}}
                \texttt{ДВИГАТЬСЯ: C, (7, 4)} \\
            }
        }
    }
    & & \\ \\

    \theutterance \stepcounter{utterance}  
    & & & \multicolumn{2}{p{0.3\linewidth}}{
        \cellcolor[rgb]{0.9,0.9,0.9}{
            \makecell[{{p{\linewidth}}}]{
                \texttt{\tiny{[GM$|$GM]}}
                \texttt{C moved to (7, 4)} \\
            }
        }
    }
    & & \\ \\

    \theutterance \stepcounter{utterance}  
    & & \multicolumn{4}{p{0.6\linewidth}}{
        \cellcolor[rgb]{0.9,0.9,0.9}{
            \makecell[{{p{\linewidth}}}]{
                \texttt{\tiny{[P1$\langle$GM]}}
                \texttt{Ваше сообщение было передано другому игроку.} \\
\\ 
\\ 
\texttt{Сейчас вы играете раунд 4 из максимум 20.} \\
\texttt{Вы вместе накопили 5 из 12 штрафов.} \\
\texttt{Другой игрок переместил объект на своей сетке.} \\
\\ 
\\ 
\texttt{Какая твоя следующая команда?} \\
            }
        }
    }
    & \\ \\

    \theutterance \stepcounter{utterance}  
    & \multicolumn{4}{p{0.6\linewidth}}{
        \cellcolor[rgb]{0.9,0.9,0.9}{
            \makecell[{{p{\linewidth}}}]{
                \texttt{\tiny{[P1$\rangle$GM]}}
                \texttt{ДВИГАТЬСЯ: C, (4, 4)} \\
            }
        }
    }
    & & \\ \\

    \theutterance \stepcounter{utterance}  
    & & & \multicolumn{2}{p{0.3\linewidth}}{
        \cellcolor[rgb]{0.9,0.9,0.9}{
            \makecell[{{p{\linewidth}}}]{
                \texttt{\tiny{[GM$|$GM]}}
                \texttt{ДВИГАТЬСЯ: C, (4, 4)} \\
            }
        }
    }
    & & \\ \\

    \theutterance \stepcounter{utterance}  
    & & & \multicolumn{2}{p{0.3\linewidth}}{
        \cellcolor[rgb]{0.9,0.9,0.9}{
            \makecell[{{p{\linewidth}}}]{
                \texttt{\tiny{[GM$|$GM]}}
                \texttt{Штраф: (4,4) не пусто, а содержит "│".} \\
\\ 
\texttt{Вы вместе накопили 6 из 12 штрафов.} \\
\texttt{Пожалуйста, попробуйте ещё раз!} \\
            }
        }
    }
    & & \\ \\

    \theutterance \stepcounter{utterance}  
    & & \multicolumn{4}{p{0.6\linewidth}}{
        \cellcolor[rgb]{0.9,0.9,0.9}{
            \makecell[{{p{\linewidth}}}]{
                \texttt{\tiny{[P1$\langle$GM]}}
                \texttt{Штраф: (4,4) не пусто, а содержит "│".} \\
\\ 
\texttt{Вы вместе накопили 6 из 12 штрафов.} \\
\texttt{Пожалуйста, попробуйте ещё раз!} \\
            }
        }
    }
    & \\ \\

    \theutterance \stepcounter{utterance}  
    & \multicolumn{4}{p{0.6\linewidth}}{
        \cellcolor[rgb]{0.9,0.9,0.9}{
            \makecell[{{p{\linewidth}}}]{
                \texttt{\tiny{[P1$\rangle$GM]}}
                \texttt{СКАЗАТЬ: Прости, я снова ошибся с проверкой линий. (4,4) занята вертикальной линией. Давай выберем полностью свободную строку. Посмотрю на 7{-}ю строку: (1,7), (2,7), (3,7), (4,7), (5,7), (6,7), (7,7). Там только '◌'. Предлагаю: H(1,7), C(2,7), T(3,7), W(4,7), I(5,7). Что скажешь?} \\
            }
        }
    }
    & & \\ \\

    \theutterance \stepcounter{utterance}  
    & & & \multicolumn{2}{p{0.3\linewidth}}{
        \cellcolor[rgb]{0.9,0.9,0.9}{
            \makecell[{{p{\linewidth}}}]{
                \texttt{\tiny{[GM$|$GM]}}
                \texttt{СКАЗАТЬ: Прости, я снова ошибся с проверкой линий. (4,4) занята вертикальной линией. Давай выберем полностью свободную строку. Посмотрю на 7{-}ю строку: (1,7), (2,7), (3,7), (4,7), (5,7), (6,7), (7,7). Там только '◌'. Предлагаю: H(1,7), C(2,7), T(3,7), W(4,7), I(5,7). Что скажешь?} \\
            }
        }
    }
    & & \\ \\

    \theutterance \stepcounter{utterance}  
    & & \multicolumn{4}{p{0.6\linewidth}}{
        \cellcolor[rgb]{0.9,0.9,0.9}{
            \makecell[{{p{\linewidth}}}]{
                \texttt{\tiny{[GM$\rangle$P2]}}
                \texttt{Успешно перемещён объект "C" в (7,4). Ваша обновлённая сетка выглядит так:} \\
\texttt{```} \\
\texttt{ 1234567} \\
\texttt{╔═╤═╤═══╗} \\
\texttt{║H│◌│◌◌◌║ 1} \\
\texttt{║◌├─┤◌◌◌║ 2} \\
\texttt{╟─┤◌├─┬─╢ 3} \\
\texttt{║◌│◌│◌│C║ 4} \\
\texttt{╟─┼─┘◌│◌║ 5} \\
\texttt{║◌│◌◌┌┴─╢ 6} \\
\texttt{║T│WI│◌◌║ 7} \\
\texttt{╚═╧══╧══╝} \\
\texttt{```} \\
\texttt{Сейчас вы играете раунд 5 из максимум 20.} \\
\texttt{Вы вместе накопили 6 из 12 штрафов.} \\
\texttt{Другой игрок отправил это сообщение:} \\
\\ 
\\ 
\texttt{"Прости, я снова ошибся с проверкой линий. (4,4) занята вертикальной линией. Давай выберем полностью свободную строку. Посмотрю на 7{-}ю строку: (1,7), (2,7), (3,7), (4,7), (5,7), (6,7), (7,7). Там только '◌'. Предлагаю: H(1,7), C(2,7), T(3,7), W(4,7), I(5,7). Что скажешь?"} \\
\\ 
\\ 
\texttt{Какая твоя следующая команда?} \\
            }
        }
    }
    & \\ \\

    \theutterance \stepcounter{utterance}  
    & & & \multicolumn{4}{p{0.6\linewidth}}{
        \cellcolor[rgb]{0.9,0.9,0.9}{
            \makecell[{{p{\linewidth}}}]{
                \texttt{\tiny{[GM$\langle$P2]}}
                \texttt{СКАЗАТЬ: Хорошая идея, 7{-}я строка выглядит свободной. Давай использовать твой план: H(1,7), C(2,7), T(3,7), W(4,7), I(5,7). Я начну перемещать свои объекты. Сначала перемещу H на (1,7).} \\
            }
        }
    }
     \\ \\

    \theutterance \stepcounter{utterance}  
    & & & \multicolumn{2}{p{0.3\linewidth}}{
        \cellcolor[rgb]{0.9,0.9,0.9}{
            \makecell[{{p{\linewidth}}}]{
                \texttt{\tiny{[GM$|$GM]}}
                \texttt{СКАЗАТЬ: Хорошая идея, 7{-}я строка выглядит свободной. Давай использовать твой план: H(1,7), C(2,7), T(3,7), W(4,7), I(5,7). Я начну перемещать свои объекты. Сначала перемещу H на (1,7).} \\
            }
        }
    }
    & & \\ \\

    \theutterance \stepcounter{utterance}  
    & & \multicolumn{4}{p{0.6\linewidth}}{
        \cellcolor[rgb]{0.9,0.9,0.9}{
            \makecell[{{p{\linewidth}}}]{
                \texttt{\tiny{[P1$\langle$GM]}}
                \texttt{Ваше сообщение было передано другому игроку.} \\
\\ 
\\ 
\texttt{Сейчас вы играете раунд 5 из максимум 20.} \\
\texttt{Вы вместе накопили 6 из 12 штрафов.} \\
\texttt{Другой игрок отправил это сообщение:} \\
\\ 
\\ 
\texttt{"Хорошая идея, 7{-}я строка выглядит свободной. Давай использовать твой план: H(1,7), C(2,7), T(3,7), W(4,7), I(5,7). Я начну перемещать свои объекты. Сначала перемещу H на (1,7)."} \\
\\ 
\\ 
\texttt{Какая твоя следующая команда?} \\
            }
        }
    }
    & \\ \\

    \theutterance \stepcounter{utterance}  
    & \multicolumn{4}{p{0.6\linewidth}}{
        \cellcolor[rgb]{0.9,0.9,0.9}{
            \makecell[{{p{\linewidth}}}]{
                \texttt{\tiny{[P1$\rangle$GM]}}
                \texttt{ДВИГАТЬСЯ: C, (2, 7)} \\
            }
        }
    }
    & & \\ \\

    \theutterance \stepcounter{utterance}  
    & & & \multicolumn{2}{p{0.3\linewidth}}{
        \cellcolor[rgb]{0.9,0.9,0.9}{
            \makecell[{{p{\linewidth}}}]{
                \texttt{\tiny{[GM$|$GM]}}
                \texttt{ДВИГАТЬСЯ: C, (2, 7)} \\
            }
        }
    }
    & & \\ \\

    \theutterance \stepcounter{utterance}  
    & & & \multicolumn{2}{p{0.3\linewidth}}{
        \cellcolor[rgb]{0.9,0.9,0.9}{
            \makecell[{{p{\linewidth}}}]{
                \texttt{\tiny{[GM$|$GM]}}
                \texttt{Штраф: (2,7) не пусто, а содержит "│".} \\
\\ 
\texttt{Вы вместе накопили 7 из 12 штрафов.} \\
\texttt{Пожалуйста, попробуйте ещё раз!} \\
            }
        }
    }
    & & \\ \\

    \theutterance \stepcounter{utterance}  
    & & \multicolumn{4}{p{0.6\linewidth}}{
        \cellcolor[rgb]{0.9,0.9,0.9}{
            \makecell[{{p{\linewidth}}}]{
                \texttt{\tiny{[P1$\langle$GM]}}
                \texttt{Штраф: (2,7) не пусто, а содержит "│".} \\
\\ 
\texttt{Вы вместе накопили 7 из 12 штрафов.} \\
\texttt{Пожалуйста, попробуйте ещё раз!} \\
            }
        }
    }
    & \\ \\

    \theutterance \stepcounter{utterance}  
    & \multicolumn{4}{p{0.6\linewidth}}{
        \cellcolor[rgb]{0.9,0.9,0.9}{
            \makecell[{{p{\linewidth}}}]{
                \texttt{\tiny{[P1$\rangle$GM]}}
                \texttt{СКАЗАТЬ: Прости, я постоянно ошибаюсь с линиями. Давай я внимательно перепроверю 7{-}ю строку. } \\
\texttt{(1,7) {-} пусто} \\
\texttt{(2,7) {-} пусто (там нет линии, я ошибся в прошлый раз, там только край сетки)} \\
\texttt{(3,7) {-} пусто} \\
\texttt{(4,7) {-} пусто} \\
\texttt{(5,7) {-} пусто} \\
\texttt{(6,7) {-} пусто} \\
\texttt{(7,7) {-} пусто} \\
\\ 
\texttt{Погоди, я вижу, что (2,7) — это часть нижней границы `╚═╧══╧══╝`. Значит, (2,7) занято.} \\
\texttt{Давай возьмем 4{-}ю строку, но только те клетки, где нет линий.} \\
\texttt{В 4{-}й строке:} \\
\texttt{(1,4) {-} пусто} \\
\texttt{(2,4) {-} линия} \\
\texttt{(3,4) {-} пусто} \\
\texttt{(4,4) {-} линия} \\
\texttt{(5,4) {-} пусто} \\
\texttt{(6,4) {-} пусто} \\
\texttt{(7,4) {-} пусто} \\
\\ 
\texttt{Это слишком сложно. Давай попробуем 7{-}ю строку, но сдвинемся вправо, чтобы не задевать границы.} \\
\texttt{Предлагаю: H(2,7) {-} нет, там граница.} \\
\texttt{Давай: H(3,7), C(4,7), T(5,7), W(6,7), I(7,7).} \\
\texttt{Проверю: (3,7), (4,7), (5,7), (6,7), (7,7) — это все '◌' в нижней части сетки. Подходит?} \\
            }
        }
    }
    & & \\ \\

    \theutterance \stepcounter{utterance}  
    & & & \multicolumn{2}{p{0.3\linewidth}}{
        \cellcolor[rgb]{0.9,0.9,0.9}{
            \makecell[{{p{\linewidth}}}]{
                \texttt{\tiny{[GM$|$GM]}}
                \texttt{СКАЗАТЬ: Прости, я постоянно ошибаюсь с линиями. Давай я внимательно перепроверю 7{-}ю строку. } \\
\texttt{(1,7) {-} пусто} \\
\texttt{(2,7) {-} пусто (там нет линии, я ошибся в прошлый раз, там только край сетки)} \\
\texttt{(3,7) {-} пусто} \\
\texttt{(4,7) {-} пусто} \\
\texttt{(5,7) {-} пусто} \\
\texttt{(6,7) {-} пусто} \\
\texttt{(7,7) {-} пусто} \\
\\ 
\texttt{Погоди, я вижу, что (2,7) — это часть нижней границы `╚═╧══╧══╝`. Значит, (2,7) занято.} \\
\texttt{Давай возьмем 4{-}ю строку, но только те клетки, где нет линий.} \\
\texttt{В 4{-}й строке:} \\
\texttt{(1,4) {-} пусто} \\
\texttt{(2,4) {-} линия} \\
\texttt{(3,4) {-} пусто} \\
\texttt{(4,4) {-} линия} \\
\texttt{(5,4) {-} пусто} \\
\texttt{(6,4) {-} пусто} \\
\texttt{(7,4) {-} пусто} \\
\\ 
\texttt{Это слишком сложно. Давай попробуем 7{-}ю строку, но сдвинемся вправо, чтобы не задевать границы.} \\
\texttt{Предлагаю: H(2,7) {-} нет, там граница.} \\
\texttt{Давай: H(3,7), C(4,7), T(5,7), W(6,7), I(7,7).} \\
\texttt{Проверю: (3,7), (4,7), (5,7), (6,7), (7,7) — это все '◌' в нижней части сетки. Подходит?} \\
            }
        }
    }
    & & \\ \\

    \theutterance \stepcounter{utterance}  
    & & & \multicolumn{2}{p{0.3\linewidth}}{
        \cellcolor[rgb]{0.9,0.9,0.9}{
            \makecell[{{p{\linewidth}}}]{
                \texttt{\tiny{[GM$|$GM]}}
                \texttt{Invalid format: tail is not empty: '} \\
\texttt{(1,7) {-} пусто} \\
\texttt{(2,7) {-} пусто (там нет линии, я ошибся в прошлый раз, там только край сетки)} \\
\texttt{(3,7) {-} пусто} \\
\texttt{(4,7) {-} пусто} \\
\texttt{(5,7) {-} пусто} \\
\texttt{(6,7) {-} пусто} \\
\texttt{(7,7) {-} пусто} \\
\\ 
\texttt{Погоди, я вижу, что (2,7) — это часть нижней границы ╚═╧══╧══╝. Значит, (2,7) занято.} \\
\texttt{Давай возьмем 4{-}ю строку, но только те клетки, где нет линий.} \\
\texttt{В 4{-}й строке:} \\
\texttt{(1,4) {-} пусто} \\
\texttt{(2,4) {-} линия} \\
\texttt{(3,4) {-} пусто} \\
\texttt{(4,4) {-} линия} \\
\texttt{(5,4) {-} пусто} \\
\texttt{(6,4) {-} пусто} \\
\texttt{(7,4) {-} пусто} \\
\\ 
\texttt{Это слишком сложно. Давай попробуем 7{-}ю строку, но сдвинемся вправо, чтобы не задевать границы.} \\
\texttt{Предлагаю: H(2,7) {-} нет, там граница.} \\
\texttt{Давай: H(3,7), C(4,7), T(5,7), W(6,7), I(7,7).} \\
\texttt{Проверю: (3,7), (4,7), (5,7), (6,7), (7,7) — это все '◌' в нижней части сетки. Подходит?} \\
            }
        }
    }
    & & \\ \\

    \theutterance \stepcounter{utterance}  
    & & \multicolumn{4}{p{0.6\linewidth}}{
        \cellcolor[rgb]{0.9,0.9,0.9}{
            \makecell[{{p{\linewidth}}}]{
                \texttt{\tiny{[P1$\langle$GM]}}
                \texttt{Штраф: Ваше сообщение не должно содержать ничего после команды!} \\
\\ 
\texttt{Убедитесь, что ваш ответ содержит только `СКАЗАТЬ: <СООБЩЕНИЕ>` или `ДВИГАТЬСЯ: <ОБЪЕКТ>, (<X>, <Y>)` и ничего больше!} \\
\\ 
\texttt{Вы вместе накопили 8 из 12 штрафов.} \\
\texttt{Пожалуйста, попробуйте ещё раз!} \\
            }
        }
    }
    & \\ \\

    \theutterance \stepcounter{utterance}  
    & \multicolumn{4}{p{0.6\linewidth}}{
        \cellcolor[rgb]{0.9,0.9,0.9}{
            \makecell[{{p{\linewidth}}}]{
                \texttt{\tiny{[P1$\rangle$GM]}}
                \texttt{СКАЗАТЬ: Прости, я запутался в линиях. Давай попробуем 4{-}ю строку, но только свободные клетки: H(1,4), T(3,4), W(5,4), I(6,4), C(7,4). Проверь, свободны ли они у тебя?} \\
            }
        }
    }
    & & \\ \\

    \theutterance \stepcounter{utterance}  
    & & & \multicolumn{2}{p{0.3\linewidth}}{
        \cellcolor[rgb]{0.9,0.9,0.9}{
            \makecell[{{p{\linewidth}}}]{
                \texttt{\tiny{[GM$|$GM]}}
                \texttt{СКАЗАТЬ: Прости, я запутался в линиях. Давай попробуем 4{-}ю строку, но только свободные клетки: H(1,4), T(3,4), W(5,4), I(6,4), C(7,4). Проверь, свободны ли они у тебя?} \\
            }
        }
    }
    & & \\ \\

    \theutterance \stepcounter{utterance}  
    & & \multicolumn{4}{p{0.6\linewidth}}{
        \cellcolor[rgb]{0.9,0.9,0.9}{
            \makecell[{{p{\linewidth}}}]{
                \texttt{\tiny{[GM$\rangle$P2]}}
                \texttt{Ваше сообщение было передано другому игроку.} \\
\\ 
\\ 
\texttt{Сейчас вы играете раунд 6 из максимум 20.} \\
\texttt{Вы вместе накопили 8 из 12 штрафов.} \\
\texttt{Другой игрок отправил это сообщение:} \\
\\ 
\\ 
\texttt{"Прости, я запутался в линиях. Давай попробуем 4{-}ю строку, но только свободные клетки: H(1,4), T(3,4), W(5,4), I(6,4), C(7,4). Проверь, свободны ли они у тебя?"} \\
\\ 
\\ 
\texttt{Какая твоя следующая команда?} \\
            }
        }
    }
    & \\ \\

    \theutterance \stepcounter{utterance}  
    & & & \multicolumn{4}{p{0.6\linewidth}}{
        \cellcolor[rgb]{0.9,0.9,0.9}{
            \makecell[{{p{\linewidth}}}]{
                \texttt{\tiny{[GM$\langle$P2]}}
                \texttt{СКАЗАТЬ: Давай проверим 4{-}ю строку у меня: (1,4) — пусто, (3,4) — пусто, (5,4) — пусто, (6,4) — пусто, (7,4) — занято моим 'C'. Так как (7,4) занято, давай просто сместим твой план на одну клетку влево или используем 7{-}ю строку, как ты предлагал раньше. На 7{-}й строке у меня: (1,7) — T, (2,7) — W, (3,7) — I, (4,7) — пусто, (5,7) — пусто, (6,7) — пусто, (7,7) — пусто. Предлагаю: H(1,7), C(2,7), T(3,7), W(4,7), I(5,7). Это кажется самым надежным.} \\
            }
        }
    }
     \\ \\

    \theutterance \stepcounter{utterance}  
    & & & \multicolumn{2}{p{0.3\linewidth}}{
        \cellcolor[rgb]{0.9,0.9,0.9}{
            \makecell[{{p{\linewidth}}}]{
                \texttt{\tiny{[GM$|$GM]}}
                \texttt{СКАЗАТЬ: Давай проверим 4{-}ю строку у меня: (1,4) — пусто, (3,4) — пусто, (5,4) — пусто, (6,4) — пусто, (7,4) — занято моим 'C'. Так как (7,4) занято, давай просто сместим твой план на одну клетку влево или используем 7{-}ю строку, как ты предлагал раньше. На 7{-}й строке у меня: (1,7) — T, (2,7) — W, (3,7) — I, (4,7) — пусто, (5,7) — пусто, (6,7) — пусто, (7,7) — пусто. Предлагаю: H(1,7), C(2,7), T(3,7), W(4,7), I(5,7). Это кажется самым надежным.} \\
            }
        }
    }
    & & \\ \\

    \theutterance \stepcounter{utterance}  
    & & \multicolumn{4}{p{0.6\linewidth}}{
        \cellcolor[rgb]{0.9,0.9,0.9}{
            \makecell[{{p{\linewidth}}}]{
                \texttt{\tiny{[P1$\langle$GM]}}
                \texttt{Ваше сообщение было передано другому игроку.} \\
\\ 
\\ 
\texttt{Сейчас вы играете раунд 6 из максимум 20.} \\
\texttt{Вы вместе накопили 8 из 12 штрафов.} \\
\texttt{Другой игрок отправил это сообщение:} \\
\\ 
\\ 
\texttt{"Давай проверим 4{-}ю строку у меня: (1,4) — пусто, (3,4) — пусто, (5,4) — пусто, (6,4) — пусто, (7,4) — занято моим 'C'. Так как (7,4) занято, давай просто сместим твой план на одну клетку влево или используем 7{-}ю строку, как ты предлагал раньше. На 7{-}й строке у меня: (1,7) — T, (2,7) — W, (3,7) — I, (4,7) — пусто, (5,7) — пусто, (6,7) — пусто, (7,7) — пусто. Предлагаю: H(1,7), C(2,7), T(3,7), W(4,7), I(5,7). Это кажется самым надежным."} \\
\\ 
\\ 
\texttt{Какая твоя следующая команда?} \\
            }
        }
    }
    & \\ \\

    \theutterance \stepcounter{utterance}  
    & \multicolumn{4}{p{0.6\linewidth}}{
        \cellcolor[rgb]{0.9,0.9,0.9}{
            \makecell[{{p{\linewidth}}}]{
                \texttt{\tiny{[P1$\rangle$GM]}}
                \texttt{СКАЗАТЬ: Понял, давай использовать 7{-}ю строку: H(1,7), C(2,7), T(3,7), W(4,7), I(5,7). Но подожди, я только что заметил, что (2,7) — это часть границы `═╧══╧══╝`. Давай выберем клетки, которые точно не задевают границы. На 7{-}й строке это (3,7), (4,7), (5,7), (6,7), (7,7). Предлагаю: H(3,7), C(4,7), T(5,7), W(6,7), I(7,7). Проверь, свободны ли они у тебя?} \\
            }
        }
    }
    & & \\ \\

    \theutterance \stepcounter{utterance}  
    & & & \multicolumn{2}{p{0.3\linewidth}}{
        \cellcolor[rgb]{0.9,0.9,0.9}{
            \makecell[{{p{\linewidth}}}]{
                \texttt{\tiny{[GM$|$GM]}}
                \texttt{СКАЗАТЬ: Понял, давай использовать 7{-}ю строку: H(1,7), C(2,7), T(3,7), W(4,7), I(5,7). Но подожди, я только что заметил, что (2,7) — это часть границы `═╧══╧══╝`. Давай выберем клетки, которые точно не задевают границы. На 7{-}й строке это (3,7), (4,7), (5,7), (6,7), (7,7). Предлагаю: H(3,7), C(4,7), T(5,7), W(6,7), I(7,7). Проверь, свободны ли они у тебя?} \\
            }
        }
    }
    & & \\ \\

    \theutterance \stepcounter{utterance}  
    & & \multicolumn{4}{p{0.6\linewidth}}{
        \cellcolor[rgb]{0.9,0.9,0.9}{
            \makecell[{{p{\linewidth}}}]{
                \texttt{\tiny{[GM$\rangle$P2]}}
                \texttt{Ваше сообщение было передано другому игроку.} \\
\\ 
\\ 
\texttt{Сейчас вы играете раунд 7 из максимум 20.} \\
\texttt{Вы вместе накопили 8 из 12 штрафов.} \\
\texttt{Другой игрок отправил это сообщение:} \\
\\ 
\\ 
\texttt{"Понял, давай использовать 7{-}ю строку: H(1,7), C(2,7), T(3,7), W(4,7), I(5,7). Но подожди, я только что заметил, что (2,7) — это часть границы ═╧══╧══╝. Давай выберем клетки, которые точно не задевают границы. На 7{-}й строке это (3,7), (4,7), (5,7), (6,7), (7,7). Предлагаю: H(3,7), C(4,7), T(5,7), W(6,7), I(7,7). Проверь, свободны ли они у тебя?"} \\
\\ 
\\ 
\texttt{Какая твоя следующая команда?} \\
            }
        }
    }
    & \\ \\

    \theutterance \stepcounter{utterance}  
    & & & \multicolumn{4}{p{0.6\linewidth}}{
        \cellcolor[rgb]{0.9,0.9,0.9}{
            \makecell[{{p{\linewidth}}}]{
                \texttt{\tiny{[GM$\langle$P2]}}
                \texttt{СКАЗАТЬ: Проверяю 7{-}ю строку у себя: (3,7) — I, (4,7) — пусто, (5,7) — пусто, (6,7) — пусто, (7,7) — пусто. (3,7) занято моим 'I'. Давай тогда возьмем клетки (4,7), (5,7), (6,7), (7,7) и добавим еще одну, например (1,7) или (2,7), если они не задевают границы. Но подожди, (1,7) и (2,7) у меня заняты 'T' и 'W'. Давай попробуем так: H(4,7), C(5,7), T(6,7), W(7,7) и I(1,7)? Нет, (1,7) занято. Давай просто: H(4,7), C(5,7), T(6,7), W(7,7) и I(3,7) — это уже почти готово! Что скажешь?} \\
            }
        }
    }
     \\ \\

    \theutterance \stepcounter{utterance}  
    & & & \multicolumn{2}{p{0.3\linewidth}}{
        \cellcolor[rgb]{0.9,0.9,0.9}{
            \makecell[{{p{\linewidth}}}]{
                \texttt{\tiny{[GM$|$GM]}}
                \texttt{СКАЗАТЬ: Проверяю 7{-}ю строку у себя: (3,7) — I, (4,7) — пусто, (5,7) — пусто, (6,7) — пусто, (7,7) — пусто. (3,7) занято моим 'I'. Давай тогда возьмем клетки (4,7), (5,7), (6,7), (7,7) и добавим еще одну, например (1,7) или (2,7), если они не задевают границы. Но подожди, (1,7) и (2,7) у меня заняты 'T' и 'W'. Давай попробуем так: H(4,7), C(5,7), T(6,7), W(7,7) и I(1,7)? Нет, (1,7) занято. Давай просто: H(4,7), C(5,7), T(6,7), W(7,7) и I(3,7) — это уже почти готово! Что скажешь?} \\
            }
        }
    }
    & & \\ \\

    \theutterance \stepcounter{utterance}  
    & & \multicolumn{4}{p{0.6\linewidth}}{
        \cellcolor[rgb]{0.9,0.9,0.9}{
            \makecell[{{p{\linewidth}}}]{
                \texttt{\tiny{[P1$\langle$GM]}}
                \texttt{Ваше сообщение было передано другому игроку.} \\
\\ 
\\ 
\texttt{Сейчас вы играете раунд 7 из максимум 20.} \\
\texttt{Вы вместе накопили 8 из 12 штрафов.} \\
\texttt{Другой игрок отправил это сообщение:} \\
\\ 
\\ 
\texttt{"Проверяю 7{-}ю строку у себя: (3,7) — I, (4,7) — пусто, (5,7) — пусто, (6,7) — пусто, (7,7) — пусто. (3,7) занято моим 'I'. Давай тогда возьмем клетки (4,7), (5,7), (6,7), (7,7) и добавим еще одну, например (1,7) или (2,7), если они не задевают границы. Но подожди, (1,7) и (2,7) у меня заняты 'T' и 'W'. Давай попробуем так: H(4,7), C(5,7), T(6,7), W(7,7) и I(1,7)? Нет, (1,7) занято. Давай просто: H(4,7), C(5,7), T(6,7), W(7,7) и I(3,7) — это уже почти готово! Что скажешь?"} \\
\\ 
\\ 
\texttt{Какая твоя следующая команда?} \\
            }
        }
    }
    & \\ \\

    \theutterance \stepcounter{utterance}  
    & \multicolumn{4}{p{0.6\linewidth}}{
        \cellcolor[rgb]{0.9,0.9,0.9}{
            \makecell[{{p{\linewidth}}}]{
                \texttt{\tiny{[P1$\rangle$GM]}}
                \texttt{СКАЗАТЬ: Хорошо, давай так: H(4,7), C(5,7), T(6,7), W(7,7), I(3,7). Я начинаю перемещать свои объекты. Сначала перемещу I на (3,7).} \\
            }
        }
    }
    & & \\ \\

    \theutterance \stepcounter{utterance}  
    & & & \multicolumn{2}{p{0.3\linewidth}}{
        \cellcolor[rgb]{0.9,0.9,0.9}{
            \makecell[{{p{\linewidth}}}]{
                \texttt{\tiny{[GM$|$GM]}}
                \texttt{СКАЗАТЬ: Хорошо, давай так: H(4,7), C(5,7), T(6,7), W(7,7), I(3,7). Я начинаю перемещать свои объекты. Сначала перемещу I на (3,7).} \\
            }
        }
    }
    & & \\ \\

    \theutterance \stepcounter{utterance}  
    & & \multicolumn{4}{p{0.6\linewidth}}{
        \cellcolor[rgb]{0.9,0.9,0.9}{
            \makecell[{{p{\linewidth}}}]{
                \texttt{\tiny{[GM$\rangle$P2]}}
                \texttt{Ваше сообщение было передано другому игроку.} \\
\\ 
\\ 
\texttt{Сейчас вы играете раунд 8 из максимум 20.} \\
\texttt{Вы вместе накопили 8 из 12 штрафов.} \\
\texttt{Другой игрок отправил это сообщение:} \\
\\ 
\\ 
\texttt{"Хорошо, давай так: H(4,7), C(5,7), T(6,7), W(7,7), I(3,7). Я начинаю перемещать свои объекты. Сначала перемещу I на (3,7)."} \\
\\ 
\\ 
\texttt{Какая твоя следующая команда?} \\
            }
        }
    }
    & \\ \\

    \theutterance \stepcounter{utterance}  
    & & & \multicolumn{4}{p{0.6\linewidth}}{
        \cellcolor[rgb]{0.9,0.9,0.9}{
            \makecell[{{p{\linewidth}}}]{
                \texttt{\tiny{[GM$\langle$P2]}}
                \texttt{ДВИГАТЬСЯ: I, (3, 7)} \\
            }
        }
    }
     \\ \\

    \theutterance \stepcounter{utterance}  
    & & & \multicolumn{2}{p{0.3\linewidth}}{
        \cellcolor[rgb]{0.9,0.9,0.9}{
            \makecell[{{p{\linewidth}}}]{
                \texttt{\tiny{[GM$|$GM]}}
                \texttt{ДВИГАТЬСЯ: I, (3, 7)} \\
            }
        }
    }
    & & \\ \\

    \theutterance \stepcounter{utterance}  
    & & & \multicolumn{2}{p{0.3\linewidth}}{
        \cellcolor[rgb]{0.9,0.9,0.9}{
            \makecell[{{p{\linewidth}}}]{
                \texttt{\tiny{[GM$|$GM]}}
                \texttt{Штраф: (3,7) не пусто, а содержит "W".} \\
\\ 
\texttt{Вы вместе накопили 9 из 12 штрафов.} \\
\texttt{Пожалуйста, попробуйте ещё раз!} \\
            }
        }
    }
    & & \\ \\

    \theutterance \stepcounter{utterance}  
    & & \multicolumn{4}{p{0.6\linewidth}}{
        \cellcolor[rgb]{0.9,0.9,0.9}{
            \makecell[{{p{\linewidth}}}]{
                \texttt{\tiny{[GM$\rangle$P2]}}
                \texttt{Штраф: (3,7) не пусто, а содержит "W".} \\
\\ 
\texttt{Вы вместе накопили 9 из 12 штрафов.} \\
\texttt{Пожалуйста, попробуйте ещё раз!} \\
            }
        }
    }
    & \\ \\

    \theutterance \stepcounter{utterance}  
    & & & \multicolumn{4}{p{0.6\linewidth}}{
        \cellcolor[rgb]{0.9,0.9,0.9}{
            \makecell[{{p{\linewidth}}}]{
                \texttt{\tiny{[GM$\langle$P2]}}
                \texttt{ДВИГАТЬСЯ: I, (4, 7)} \\
            }
        }
    }
     \\ \\

    \theutterance \stepcounter{utterance}  
    & & & \multicolumn{2}{p{0.3\linewidth}}{
        \cellcolor[rgb]{0.9,0.9,0.9}{
            \makecell[{{p{\linewidth}}}]{
                \texttt{\tiny{[GM$|$GM]}}
                \texttt{ДВИГАТЬСЯ: I, (4, 7)} \\
            }
        }
    }
    & & \\ \\

    \theutterance \stepcounter{utterance}  
    & & & \multicolumn{2}{p{0.3\linewidth}}{
        \cellcolor[rgb]{0.9,0.9,0.9}{
            \makecell[{{p{\linewidth}}}]{
                \texttt{\tiny{[GM$|$GM]}}
                \texttt{Штраф: (4,7) не пусто, а содержит "I".} \\
\\ 
\texttt{Вы вместе накопили 10 из 12 штрафов.} \\
\texttt{Пожалуйста, попробуйте ещё раз!} \\
            }
        }
    }
    & & \\ \\

    \theutterance \stepcounter{utterance}  
    & & \multicolumn{4}{p{0.6\linewidth}}{
        \cellcolor[rgb]{0.9,0.9,0.9}{
            \makecell[{{p{\linewidth}}}]{
                \texttt{\tiny{[GM$\rangle$P2]}}
                \texttt{Штраф: (4,7) не пусто, а содержит "I".} \\
\\ 
\texttt{Вы вместе накопили 10 из 12 штрафов.} \\
\texttt{Пожалуйста, попробуйте ещё раз!} \\
            }
        }
    }
    & \\ \\

    \theutterance \stepcounter{utterance}  
    & & & \multicolumn{4}{p{0.6\linewidth}}{
        \cellcolor[rgb]{0.9,0.9,0.9}{
            \makecell[{{p{\linewidth}}}]{
                \texttt{\tiny{[GM$\langle$P2]}}
                \texttt{ДВИГАТЬСЯ: I, (3, 7)} \\
\\ 
\texttt{(Примечание: Я пытаюсь переместить I в (3,7), но на моей сетке в (3,7) находится 'I'. По правилам, я должен переместить объект в пустую клетку. Посмотрю на сетку: (3,7) — это 'I', (4,7) — это 'W', (5,7) — это 'I'. Ой, я перепутал координаты. } \\
\\ 
\texttt{Моя сетка:} \\
\texttt{(1,7) = T} \\
\texttt{(2,7) = W} \\
\texttt{(3,7) = I} \\
\texttt{(4,7) = пусто} \\
\texttt{(5,7) = пусто} \\
\texttt{(6,7) = пусто} \\
\texttt{(7,7) = пусто} \\
\\ 
\texttt{Цель: I(3,7). Но I уже на (3,7). Значит, мне нужно переместить другие объекты.} \\
\texttt{Сначала перемещу H(1,1) на (4,7).)} \\
\\ 
\texttt{ДВИГАТЬСЯ: H, (4, 7)} \\
            }
        }
    }
     \\ \\

    \theutterance \stepcounter{utterance}  
    & & & \multicolumn{2}{p{0.3\linewidth}}{
        \cellcolor[rgb]{0.9,0.9,0.9}{
            \makecell[{{p{\linewidth}}}]{
                \texttt{\tiny{[GM$|$GM]}}
                \texttt{ДВИГАТЬСЯ: I, (3, 7)} \\
\\ 
\texttt{(Примечание: Я пытаюсь переместить I в (3,7), но на моей сетке в (3,7) находится 'I'. По правилам, я должен переместить объект в пустую клетку. Посмотрю на сетку: (3,7) — это 'I', (4,7) — это 'W', (5,7) — это 'I'. Ой, я перепутал координаты. } \\
\\ 
\texttt{Моя сетка:} \\
\texttt{(1,7) = T} \\
\texttt{(2,7) = W} \\
\texttt{(3,7) = I} \\
\texttt{(4,7) = пусто} \\
\texttt{(5,7) = пусто} \\
\texttt{(6,7) = пусто} \\
\texttt{(7,7) = пусто} \\
\\ 
\texttt{Цель: I(3,7). Но I уже на (3,7). Значит, мне нужно переместить другие объекты.} \\
\texttt{Сначала перемещу H(1,1) на (4,7).)} \\
\\ 
\texttt{ДВИГАТЬСЯ: H, (4, 7)} \\
            }
        }
    }
    & & \\ \\

    \theutterance \stepcounter{utterance}  
    & & & \multicolumn{2}{p{0.3\linewidth}}{
        \cellcolor[rgb]{0.9,0.9,0.9}{
            \makecell[{{p{\linewidth}}}]{
                \texttt{\tiny{[GM$|$GM]}}
                \texttt{Invalid format: head is not empty: 'ДВИГАТЬСЯ: I, (3, 7)} \\
\\ 
\texttt{(Примечание: Я пытаюсь переместить I в (3,7), но на моей сетке в (3,7) находится 'I'. По правилам, я должен переместить объект в пустую клетку. Посмотрю на сетку: (3,7) — это 'I', (4,7) — это 'W', (5,7) — это 'I'. Ой, я перепутал координаты. } \\
\\ 
\texttt{Моя сетка:} \\
\texttt{(1,7) = T} \\
\texttt{(2,7) = W} \\
\texttt{(3,7) = I} \\
\texttt{(4,7) = пусто} \\
\texttt{(5,7) = пусто} \\
\texttt{(6,7) = пусто} \\
\texttt{(7,7) = пусто} \\
\\ 
\texttt{Цель: I(3,7). Но I уже на (3,7). Значит, мне нужно переместить другие объекты.} \\
\texttt{Сначала перемещу H(1,1) на (4,7).)} \\
\\ 
\texttt{'} \\
            }
        }
    }
    & & \\ \\

    \theutterance \stepcounter{utterance}  
    & & \multicolumn{4}{p{0.6\linewidth}}{
        \cellcolor[rgb]{0.9,0.9,0.9}{
            \makecell[{{p{\linewidth}}}]{
                \texttt{\tiny{[GM$\rangle$P2]}}
                \texttt{Штраф: Ваше сообщение не должно содержать ничего перед командой!} \\
\\ 
\texttt{Убедитесь, что ваш ответ содержит только `СКАЗАТЬ: <СООБЩЕНИЕ>` или `ДВИГАТЬСЯ: <ОБЪЕКТ>, (<X>, <Y>)` и ничего больше!} \\
\\ 
\texttt{Вы вместе накопили 11 из 12 штрафов.} \\
\texttt{Пожалуйста, попробуйте ещё раз!} \\
            }
        }
    }
    & \\ \\

    \theutterance \stepcounter{utterance}  
    & & & \multicolumn{4}{p{0.6\linewidth}}{
        \cellcolor[rgb]{0.9,0.9,0.9}{
            \makecell[{{p{\linewidth}}}]{
                \texttt{\tiny{[GM$\langle$P2]}}
                \texttt{ДВИГАТЬСЯ: H, (4, 7)} \\
            }
        }
    }
     \\ \\

    \theutterance \stepcounter{utterance}  
    & & & \multicolumn{2}{p{0.3\linewidth}}{
        \cellcolor[rgb]{0.9,0.9,0.9}{
            \makecell[{{p{\linewidth}}}]{
                \texttt{\tiny{[GM$|$GM]}}
                \texttt{ДВИГАТЬСЯ: H, (4, 7)} \\
            }
        }
    }
    & & \\ \\

    \theutterance \stepcounter{utterance}  
    & & & \multicolumn{2}{p{0.3\linewidth}}{
        \cellcolor[rgb]{0.9,0.9,0.9}{
            \makecell[{{p{\linewidth}}}]{
                \texttt{\tiny{[GM$|$GM]}}
                \texttt{Штраф: (4,7) не пусто, а содержит "I".} \\
\\ 
\texttt{Вы вместе накопили 12 из 12 штрафов.} \\
\texttt{Пожалуйста, попробуйте ещё раз!} \\
            }
        }
    }
    & & \\ \\

    \theutterance \stepcounter{utterance}  
    & & \multicolumn{4}{p{0.6\linewidth}}{
        \cellcolor[rgb]{0.9,0.9,0.9}{
            \makecell[{{p{\linewidth}}}]{
                \texttt{\tiny{[GM$\rangle$P2]}}
                \texttt{Штраф: (4,7) не пусто, а содержит "I".} \\
\\ 
\texttt{Вы вместе накопили 12 из 12 штрафов.} \\
\texttt{Пожалуйста, попробуйте ещё раз!} \\
            }
        }
    }
    & \\ \\

    \theutterance \stepcounter{utterance}  
    & & & \multicolumn{4}{p{0.6\linewidth}}{
        \cellcolor[rgb]{0.9,0.9,0.9}{
            \makecell[{{p{\linewidth}}}]{
                \texttt{\tiny{[GM$\langle$P2]}}
                \texttt{ДВИГАТЬСЯ: H, (5, 7)} \\
            }
        }
    }
     \\ \\

    \theutterance \stepcounter{utterance}  
    & & & \multicolumn{2}{p{0.3\linewidth}}{
        \cellcolor[rgb]{0.9,0.9,0.9}{
            \makecell[{{p{\linewidth}}}]{
                \texttt{\tiny{[GM$|$GM]}}
                \texttt{ДВИГАТЬСЯ: H, (5, 7)} \\
            }
        }
    }
    & & \\ \\

    \theutterance \stepcounter{utterance}  
    & & & \multicolumn{2}{p{0.3\linewidth}}{
        \cellcolor[rgb]{0.9,0.9,0.9}{
            \makecell[{{p{\linewidth}}}]{
                \texttt{\tiny{[GM$|$GM]}}
                \texttt{Штраф: (5,7) не пусто, а содержит "│".} \\
\\ 
\texttt{Вы вместе накопили 13 из 12 штрафов.} \\
\texttt{Пожалуйста, попробуйте ещё раз!} \\
            }
        }
    }
    & & \\ \\

    \theutterance \stepcounter{utterance}  
    & & \multicolumn{4}{p{0.6\linewidth}}{
        \cellcolor[rgb]{0.9,0.9,0.9}{
            \makecell[{{p{\linewidth}}}]{
                \texttt{\tiny{[GM$\rangle$P2]}}
                \texttt{Штраф: (5,7) не пусто, а содержит "│".} \\
\\ 
\texttt{Вы вместе накопили 13 из 12 штрафов.} \\
\texttt{Пожалуйста, попробуйте ещё раз!} \\
            }
        }
    }
    & \\ \\

    \theutterance \stepcounter{utterance}  
    & & & \multicolumn{4}{p{0.6\linewidth}}{
        \cellcolor[rgb]{0.9,0.9,0.9}{
            \makecell[{{p{\linewidth}}}]{
                \texttt{\tiny{[GM$\langle$P2]}}
                \texttt{ДВИГАТЬСЯ: H, (6, 7)} \\
            }
        }
    }
     \\ \\

    \theutterance \stepcounter{utterance}  
    & & & \multicolumn{2}{p{0.3\linewidth}}{
        \cellcolor[rgb]{0.9,0.9,0.9}{
            \makecell[{{p{\linewidth}}}]{
                \texttt{\tiny{[GM$|$GM]}}
                \texttt{ДВИГАТЬСЯ: H, (6, 7)} \\
            }
        }
    }
    & & \\ \\

    \theutterance \stepcounter{utterance}  
    & & & \multicolumn{2}{p{0.3\linewidth}}{
        \cellcolor[rgb]{0.9,0.9,0.9}{
            \makecell[{{p{\linewidth}}}]{
                \texttt{\tiny{[GM$|$GM]}}
                \texttt{H moved to (6, 7)} \\
            }
        }
    }
    & & \\ \\

    \theutterance \stepcounter{utterance}  
    & & \multicolumn{4}{p{0.6\linewidth}}{
        \cellcolor[rgb]{0.9,0.9,0.9}{
            \makecell[{{p{\linewidth}}}]{
                \texttt{\tiny{[P1$\langle$GM]}}
                \texttt{Ваше сообщение было передано другому игроку.} \\
\\ 
\\ 
\texttt{Сейчас вы играете раунд 8 из максимум 20.} \\
\texttt{Вы вместе накопили 13 из 12 штрафов.} \\
\texttt{Другой игрок переместил объект на своей сетке.} \\
\\ 
\\ 
\texttt{Какая твоя следующая команда?} \\
            }
        }
    }
    & \\ \\

    \theutterance \stepcounter{utterance}  
    & \multicolumn{4}{p{0.6\linewidth}}{
        \cellcolor[rgb]{0.9,0.9,0.9}{
            \makecell[{{p{\linewidth}}}]{
                \texttt{\tiny{[P1$\rangle$GM]}}
                \texttt{ДВИГАТЬСЯ: I, (3, 7)} \\
            }
        }
    }
    & & \\ \\

    \theutterance \stepcounter{utterance}  
    & & & \multicolumn{2}{p{0.3\linewidth}}{
        \cellcolor[rgb]{0.9,0.9,0.9}{
            \makecell[{{p{\linewidth}}}]{
                \texttt{\tiny{[GM$|$GM]}}
                \texttt{ДВИГАТЬСЯ: I, (3, 7)} \\
            }
        }
    }
    & & \\ \\

    \theutterance \stepcounter{utterance}  
    & & & \multicolumn{2}{p{0.3\linewidth}}{
        \cellcolor[rgb]{0.9,0.9,0.9}{
            \makecell[{{p{\linewidth}}}]{
                \texttt{\tiny{[GM$|$GM]}}
                \texttt{I moved to (3, 7)} \\
            }
        }
    }
    & & \\ \\

    \theutterance \stepcounter{utterance}  
    & & \multicolumn{4}{p{0.6\linewidth}}{
        \cellcolor[rgb]{0.9,0.9,0.9}{
            \makecell[{{p{\linewidth}}}]{
                \texttt{\tiny{[GM$\rangle$P2]}}
                \texttt{Успешно перемещён объект "H" в (6,7). Ваша обновлённая сетка выглядит так:} \\
\texttt{```} \\
\texttt{ 1234567} \\
\texttt{╔═╤═╤═══╗} \\
\texttt{║◌│◌│◌◌◌║ 1} \\
\texttt{║◌├─┤◌◌◌║ 2} \\
\texttt{╟─┤◌├─┬─╢ 3} \\
\texttt{║◌│◌│◌│C║ 4} \\
\texttt{╟─┼─┘◌│◌║ 5} \\
\texttt{║◌│◌◌┌┴─╢ 6} \\
\texttt{║T│WI│H◌║ 7} \\
\texttt{╚═╧══╧══╝} \\
\texttt{```} \\
\texttt{Сейчас вы играете раунд 9 из максимум 20.} \\
\texttt{Вы вместе накопили 13 из 12 штрафов.} \\
\texttt{Другой игрок переместил объект на своей сетке.} \\
\\ 
\\ 
\texttt{Какая твоя следующая команда?} \\
            }
        }
    }
    & \\ \\

    \theutterance \stepcounter{utterance}  
    & & & \multicolumn{4}{p{0.6\linewidth}}{
        \cellcolor[rgb]{0.9,0.9,0.9}{
            \makecell[{{p{\linewidth}}}]{
                \texttt{\tiny{[GM$\langle$P2]}}
                \texttt{СКАЗАТЬ: Мы превысили лимит штрафов, но давай попробуем закончить, если сможем. Мои текущие позиции: T(1,7), W(2,7), I(3,7), H(6,7), C(7,4). Наша цель была: H(4,7), C(5,7), T(6,7), W(7,7), I(3,7). Я буду двигать объекты к этим целям. Сначала попробую переместить T на (6,7), но там уже H. Давай сначала перемещу W на (7,7).} \\
            }
        }
    }
     \\ \\

    \theutterance \stepcounter{utterance}  
    & & & \multicolumn{2}{p{0.3\linewidth}}{
        \cellcolor[rgb]{0.9,0.9,0.9}{
            \makecell[{{p{\linewidth}}}]{
                \texttt{\tiny{[GM$|$GM]}}
                \texttt{СКАЗАТЬ: Мы превысили лимит штрафов, но давай попробуем закончить, если сможем. Мои текущие позиции: T(1,7), W(2,7), I(3,7), H(6,7), C(7,4). Наша цель была: H(4,7), C(5,7), T(6,7), W(7,7), I(3,7). Я буду двигать объекты к этим целям. Сначала попробую переместить T на (6,7), но там уже H. Давай сначала перемещу W на (7,7).} \\
            }
        }
    }
    & & \\ \\

    \theutterance \stepcounter{utterance}  
    & & \multicolumn{4}{p{0.6\linewidth}}{
        \cellcolor[rgb]{0.9,0.9,0.9}{
            \makecell[{{p{\linewidth}}}]{
                \texttt{\tiny{[P1$\langle$GM]}}
                \texttt{Успешно перемещён объект "I" в (3,7). Ваша обновлённая сетка выглядит так:} \\
\texttt{```} \\
\texttt{ 1234567} \\
\texttt{╔═╤═╤═══╗} \\
\texttt{║C│◌│◌T◌║ 1} \\
\texttt{║◌├─┤W◌◌║ 2} \\
\texttt{╟─┤◌├─┬─╢ 3} \\
\texttt{║◌│◌│◌│◌║ 4} \\
\texttt{╟─┼─┘◌│◌║ 5} \\
\texttt{║◌│◌H┌┴─╢ 6} \\
\texttt{║◌│I◌│◌◌║ 7} \\
\texttt{╚═╧══╧══╝} \\
\texttt{```} \\
\texttt{Сейчас вы играете раунд 9 из максимум 20.} \\
\texttt{Вы вместе накопили 13 из 12 штрафов.} \\
\texttt{Другой игрок отправил это сообщение:} \\
\\ 
\\ 
\texttt{"Мы превысили лимит штрафов, но давай попробуем закончить, если сможем. Мои текущие позиции: T(1,7), W(2,7), I(3,7), H(6,7), C(7,4). Наша цель была: H(4,7), C(5,7), T(6,7), W(7,7), I(3,7). Я буду двигать объекты к этим целям. Сначала попробую переместить T на (6,7), но там уже H. Давай сначала перемещу W на (7,7)."} \\
\\ 
\\ 
\texttt{Какая твоя следующая команда?} \\
            }
        }
    }
    & \\ \\

    \theutterance \stepcounter{utterance}  
    & \multicolumn{4}{p{0.6\linewidth}}{
        \cellcolor[rgb]{0.9,0.9,0.9}{
            \makecell[{{p{\linewidth}}}]{
                \texttt{\tiny{[P1$\rangle$GM]}}
                \texttt{ДВИГАТЬСЯ: I, (3, 7)} \\
            }
        }
    }
    & & \\ \\

    \theutterance \stepcounter{utterance}  
    & & & \multicolumn{2}{p{0.3\linewidth}}{
        \cellcolor[rgb]{0.9,0.9,0.9}{
            \makecell[{{p{\linewidth}}}]{
                \texttt{\tiny{[GM$|$GM]}}
                \texttt{ДВИГАТЬСЯ: I, (3, 7)} \\
            }
        }
    }
    & & \\ \\

    \theutterance \stepcounter{utterance}  
    & & & \multicolumn{2}{p{0.3\linewidth}}{
        \cellcolor[rgb]{0.9,0.9,0.9}{
            \makecell[{{p{\linewidth}}}]{
                \texttt{\tiny{[GM$|$GM]}}
                \texttt{Штраф: (3,7) не пусто, а содержит "I".} \\
\\ 
\texttt{Вы вместе накопили 14 из 12 штрафов.} \\
\texttt{Пожалуйста, попробуйте ещё раз!} \\
            }
        }
    }
    & & \\ \\

    \theutterance \stepcounter{utterance}  
    & & \multicolumn{4}{p{0.6\linewidth}}{
        \cellcolor[rgb]{0.9,0.9,0.9}{
            \makecell[{{p{\linewidth}}}]{
                \texttt{\tiny{[P1$\langle$GM]}}
                \texttt{Штраф: (3,7) не пусто, а содержит "I".} \\
\\ 
\texttt{Вы вместе накопили 14 из 12 штрафов.} \\
\texttt{Пожалуйста, попробуйте ещё раз!} \\
            }
        }
    }
    & \\ \\

    \theutterance \stepcounter{utterance}  
    & \multicolumn{4}{p{0.6\linewidth}}{
        \cellcolor[rgb]{0.9,0.9,0.9}{
            \makecell[{{p{\linewidth}}}]{
                \texttt{\tiny{[P1$\rangle$GM]}}
                \texttt{ДВИГАТЬСЯ: T, (6, 7)} \\
            }
        }
    }
    & & \\ \\

    \theutterance \stepcounter{utterance}  
    & & & \multicolumn{2}{p{0.3\linewidth}}{
        \cellcolor[rgb]{0.9,0.9,0.9}{
            \makecell[{{p{\linewidth}}}]{
                \texttt{\tiny{[GM$|$GM]}}
                \texttt{ДВИГАТЬСЯ: T, (6, 7)} \\
            }
        }
    }
    & & \\ \\

    \theutterance \stepcounter{utterance}  
    & & & \multicolumn{2}{p{0.3\linewidth}}{
        \cellcolor[rgb]{0.9,0.9,0.9}{
            \makecell[{{p{\linewidth}}}]{
                \texttt{\tiny{[GM$|$GM]}}
                \texttt{T moved to (6, 7)} \\
            }
        }
    }
    & & \\ \\

    \theutterance \stepcounter{utterance}  
    & & \multicolumn{4}{p{0.6\linewidth}}{
        \cellcolor[rgb]{0.9,0.9,0.9}{
            \makecell[{{p{\linewidth}}}]{
                \texttt{\tiny{[GM$\rangle$P2]}}
                \texttt{Ваше сообщение было передано другому игроку.} \\
\\ 
\\ 
\texttt{Сейчас вы играете раунд 10 из максимум 20.} \\
\texttt{Вы вместе накопили 14 из 12 штрафов.} \\
\texttt{Другой игрок переместил объект на своей сетке.} \\
\\ 
\\ 
\texttt{Какая твоя следующая команда?} \\
            }
        }
    }
    & \\ \\

    \theutterance \stepcounter{utterance}  
    & & & \multicolumn{4}{p{0.6\linewidth}}{
        \cellcolor[rgb]{0.9,0.9,0.9}{
            \makecell[{{p{\linewidth}}}]{
                \texttt{\tiny{[GM$\langle$P2]}}
                \texttt{ДВИГАТЬСЯ: W, (7, 7)} \\
            }
        }
    }
     \\ \\

    \theutterance \stepcounter{utterance}  
    & & & \multicolumn{2}{p{0.3\linewidth}}{
        \cellcolor[rgb]{0.9,0.9,0.9}{
            \makecell[{{p{\linewidth}}}]{
                \texttt{\tiny{[GM$|$GM]}}
                \texttt{ДВИГАТЬСЯ: W, (7, 7)} \\
            }
        }
    }
    & & \\ \\

    \theutterance \stepcounter{utterance}  
    & & & \multicolumn{2}{p{0.3\linewidth}}{
        \cellcolor[rgb]{0.9,0.9,0.9}{
            \makecell[{{p{\linewidth}}}]{
                \texttt{\tiny{[GM$|$GM]}}
                \texttt{W moved to (7, 7)} \\
            }
        }
    }
    & & \\ \\

    \theutterance \stepcounter{utterance}  
    & & \multicolumn{4}{p{0.6\linewidth}}{
        \cellcolor[rgb]{0.9,0.9,0.9}{
            \makecell[{{p{\linewidth}}}]{
                \texttt{\tiny{[P1$\langle$GM]}}
                \texttt{Успешно перемещён объект "T" в (6,7). Ваша обновлённая сетка выглядит так:} \\
\texttt{```} \\
\texttt{ 1234567} \\
\texttt{╔═╤═╤═══╗} \\
\texttt{║C│◌│◌◌◌║ 1} \\
\texttt{║◌├─┤W◌◌║ 2} \\
\texttt{╟─┤◌├─┬─╢ 3} \\
\texttt{║◌│◌│◌│◌║ 4} \\
\texttt{╟─┼─┘◌│◌║ 5} \\
\texttt{║◌│◌H┌┴─╢ 6} \\
\texttt{║◌│I◌│T◌║ 7} \\
\texttt{╚═╧══╧══╝} \\
\texttt{```} \\
\texttt{Сейчас вы играете раунд 10 из максимум 20.} \\
\texttt{Вы вместе накопили 14 из 12 штрафов.} \\
\texttt{Другой игрок переместил объект на своей сетке.} \\
\\ 
\\ 
\texttt{Какая твоя следующая команда?} \\
            }
        }
    }
    & \\ \\

    \theutterance \stepcounter{utterance}  
    & \multicolumn{4}{p{0.6\linewidth}}{
        \cellcolor[rgb]{0.9,0.9,0.9}{
            \makecell[{{p{\linewidth}}}]{
                \texttt{\tiny{[P1$\rangle$GM]}}
                \texttt{ДВИГАТЬСЯ: W, (7, 7)} \\
            }
        }
    }
    & & \\ \\

    \theutterance \stepcounter{utterance}  
    & & & \multicolumn{2}{p{0.3\linewidth}}{
        \cellcolor[rgb]{0.9,0.9,0.9}{
            \makecell[{{p{\linewidth}}}]{
                \texttt{\tiny{[GM$|$GM]}}
                \texttt{ДВИГАТЬСЯ: W, (7, 7)} \\
            }
        }
    }
    & & \\ \\

    \theutterance \stepcounter{utterance}  
    & & & \multicolumn{2}{p{0.3\linewidth}}{
        \cellcolor[rgb]{0.9,0.9,0.9}{
            \makecell[{{p{\linewidth}}}]{
                \texttt{\tiny{[GM$|$GM]}}
                \texttt{W moved to (7, 7)} \\
            }
        }
    }
    & & \\ \\

    \theutterance \stepcounter{utterance}  
    & & \multicolumn{4}{p{0.6\linewidth}}{
        \cellcolor[rgb]{0.9,0.9,0.9}{
            \makecell[{{p{\linewidth}}}]{
                \texttt{\tiny{[GM$\rangle$P2]}}
                \texttt{Успешно перемещён объект "W" в (7,7). Ваша обновлённая сетка выглядит так:} \\
\texttt{```} \\
\texttt{ 1234567} \\
\texttt{╔═╤═╤═══╗} \\
\texttt{║◌│◌│◌◌◌║ 1} \\
\texttt{║◌├─┤◌◌◌║ 2} \\
\texttt{╟─┤◌├─┬─╢ 3} \\
\texttt{║◌│◌│◌│C║ 4} \\
\texttt{╟─┼─┘◌│◌║ 5} \\
\texttt{║◌│◌◌┌┴─╢ 6} \\
\texttt{║T│◌I│HW║ 7} \\
\texttt{╚═╧══╧══╝} \\
\texttt{```} \\
\texttt{Сейчас вы играете раунд 11 из максимум 20.} \\
\texttt{Вы вместе накопили 14 из 12 штрафов.} \\
\texttt{Другой игрок переместил объект на своей сетке.} \\
\\ 
\\ 
\texttt{Какая твоя следующая команда?} \\
            }
        }
    }
    & \\ \\

    \theutterance \stepcounter{utterance}  
    & & & \multicolumn{4}{p{0.6\linewidth}}{
        \cellcolor[rgb]{0.9,0.9,0.9}{
            \makecell[{{p{\linewidth}}}]{
                \texttt{\tiny{[GM$\langle$P2]}}
                \texttt{ДВИГАТЬСЯ: T, (6, 7)} \\
            }
        }
    }
     \\ \\

    \theutterance \stepcounter{utterance}  
    & & & \multicolumn{2}{p{0.3\linewidth}}{
        \cellcolor[rgb]{0.9,0.9,0.9}{
            \makecell[{{p{\linewidth}}}]{
                \texttt{\tiny{[GM$|$GM]}}
                \texttt{ДВИГАТЬСЯ: T, (6, 7)} \\
            }
        }
    }
    & & \\ \\

    \theutterance \stepcounter{utterance}  
    & & & \multicolumn{2}{p{0.3\linewidth}}{
        \cellcolor[rgb]{0.9,0.9,0.9}{
            \makecell[{{p{\linewidth}}}]{
                \texttt{\tiny{[GM$|$GM]}}
                \texttt{Штраф: (6,7) не пусто, а содержит "H".} \\
\\ 
\texttt{Вы вместе накопили 15 из 12 штрафов.} \\
\texttt{Пожалуйста, попробуйте ещё раз!} \\
            }
        }
    }
    & & \\ \\

    \theutterance \stepcounter{utterance}  
    & & \multicolumn{4}{p{0.6\linewidth}}{
        \cellcolor[rgb]{0.9,0.9,0.9}{
            \makecell[{{p{\linewidth}}}]{
                \texttt{\tiny{[GM$\rangle$P2]}}
                \texttt{Штраф: (6,7) не пусто, а содержит "H".} \\
\\ 
\texttt{Вы вместе накопили 15 из 12 штрафов.} \\
\texttt{Пожалуйста, попробуйте ещё раз!} \\
            }
        }
    }
    & \\ \\

    \theutterance \stepcounter{utterance}  
    & & & \multicolumn{4}{p{0.6\linewidth}}{
        \cellcolor[rgb]{0.9,0.9,0.9}{
            \makecell[{{p{\linewidth}}}]{
                \texttt{\tiny{[GM$\langle$P2]}}
                \texttt{ДВИГАТЬСЯ: T, (5, 7)} \\
            }
        }
    }
     \\ \\

    \theutterance \stepcounter{utterance}  
    & & & \multicolumn{2}{p{0.3\linewidth}}{
        \cellcolor[rgb]{0.9,0.9,0.9}{
            \makecell[{{p{\linewidth}}}]{
                \texttt{\tiny{[GM$|$GM]}}
                \texttt{ДВИГАТЬСЯ: T, (5, 7)} \\
            }
        }
    }
    & & \\ \\

    \theutterance \stepcounter{utterance}  
    & & & \multicolumn{2}{p{0.3\linewidth}}{
        \cellcolor[rgb]{0.9,0.9,0.9}{
            \makecell[{{p{\linewidth}}}]{
                \texttt{\tiny{[GM$|$GM]}}
                \texttt{Штраф: (5,7) не пусто, а содержит "│".} \\
\\ 
\texttt{Вы вместе накопили 16 из 12 штрафов.} \\
\texttt{Пожалуйста, попробуйте ещё раз!} \\
            }
        }
    }
    & & \\ \\

    \theutterance \stepcounter{utterance}  
    & & \multicolumn{4}{p{0.6\linewidth}}{
        \cellcolor[rgb]{0.9,0.9,0.9}{
            \makecell[{{p{\linewidth}}}]{
                \texttt{\tiny{[GM$\rangle$P2]}}
                \texttt{Штраф: (5,7) не пусто, а содержит "│".} \\
\\ 
\texttt{Вы вместе накопили 16 из 12 штрафов.} \\
\texttt{Пожалуйста, попробуйте ещё раз!} \\
            }
        }
    }
    & \\ \\

    \theutterance \stepcounter{utterance}  
    & & & \multicolumn{4}{p{0.6\linewidth}}{
        \cellcolor[rgb]{0.9,0.9,0.9}{
            \makecell[{{p{\linewidth}}}]{
                \texttt{\tiny{[GM$\langle$P2]}}
                \texttt{ДВИГАТЬСЯ: T, (4, 7)} \\
            }
        }
    }
     \\ \\

    \theutterance \stepcounter{utterance}  
    & & & \multicolumn{2}{p{0.3\linewidth}}{
        \cellcolor[rgb]{0.9,0.9,0.9}{
            \makecell[{{p{\linewidth}}}]{
                \texttt{\tiny{[GM$|$GM]}}
                \texttt{ДВИГАТЬСЯ: T, (4, 7)} \\
            }
        }
    }
    & & \\ \\

    \theutterance \stepcounter{utterance}  
    & & & \multicolumn{2}{p{0.3\linewidth}}{
        \cellcolor[rgb]{0.9,0.9,0.9}{
            \makecell[{{p{\linewidth}}}]{
                \texttt{\tiny{[GM$|$GM]}}
                \texttt{Штраф: (4,7) не пусто, а содержит "I".} \\
\\ 
\texttt{Вы вместе накопили 17 из 12 штрафов.} \\
\texttt{Пожалуйста, попробуйте ещё раз!} \\
            }
        }
    }
    & & \\ \\

    \theutterance \stepcounter{utterance}  
    & & \multicolumn{4}{p{0.6\linewidth}}{
        \cellcolor[rgb]{0.9,0.9,0.9}{
            \makecell[{{p{\linewidth}}}]{
                \texttt{\tiny{[GM$\rangle$P2]}}
                \texttt{Штраф: (4,7) не пусто, а содержит "I".} \\
\\ 
\texttt{Вы вместе накопили 17 из 12 штрафов.} \\
\texttt{Пожалуйста, попробуйте ещё раз!} \\
            }
        }
    }
    & \\ \\

    \theutterance \stepcounter{utterance}  
    & & & \multicolumn{4}{p{0.6\linewidth}}{
        \cellcolor[rgb]{0.9,0.9,0.9}{
            \makecell[{{p{\linewidth}}}]{
                \texttt{\tiny{[GM$\langle$P2]}}
                \texttt{ДВИГАТЬСЯ: T, (5, 6)} \\
            }
        }
    }
     \\ \\

    \theutterance \stepcounter{utterance}  
    & & & \multicolumn{2}{p{0.3\linewidth}}{
        \cellcolor[rgb]{0.9,0.9,0.9}{
            \makecell[{{p{\linewidth}}}]{
                \texttt{\tiny{[GM$|$GM]}}
                \texttt{ДВИГАТЬСЯ: T, (5, 6)} \\
            }
        }
    }
    & & \\ \\

    \theutterance \stepcounter{utterance}  
    & & & \multicolumn{2}{p{0.3\linewidth}}{
        \cellcolor[rgb]{0.9,0.9,0.9}{
            \makecell[{{p{\linewidth}}}]{
                \texttt{\tiny{[GM$|$GM]}}
                \texttt{Штраф: (5,6) не пусто, а содержит "┌".} \\
\\ 
\texttt{Вы вместе накопили 18 из 12 штрафов.} \\
\texttt{Пожалуйста, попробуйте ещё раз!} \\
            }
        }
    }
    & & \\ \\

    \theutterance \stepcounter{utterance}  
    & & \multicolumn{4}{p{0.6\linewidth}}{
        \cellcolor[rgb]{0.9,0.9,0.9}{
            \makecell[{{p{\linewidth}}}]{
                \texttt{\tiny{[GM$\rangle$P2]}}
                \texttt{Штраф: (5,6) не пусто, а содержит "┌".} \\
\\ 
\texttt{Вы вместе накопили 18 из 12 штрафов.} \\
\texttt{Пожалуйста, попробуйте ещё раз!} \\
            }
        }
    }
    & \\ \\

    \theutterance \stepcounter{utterance}  
    & & & \multicolumn{4}{p{0.6\linewidth}}{
        \cellcolor[rgb]{0.9,0.9,0.9}{
            \makecell[{{p{\linewidth}}}]{
                \texttt{\tiny{[GM$\langle$P2]}}
                \texttt{ДВИГАТЬСЯ: T, (1, 6)} \\
            }
        }
    }
     \\ \\

    \theutterance \stepcounter{utterance}  
    & & & \multicolumn{2}{p{0.3\linewidth}}{
        \cellcolor[rgb]{0.9,0.9,0.9}{
            \makecell[{{p{\linewidth}}}]{
                \texttt{\tiny{[GM$|$GM]}}
                \texttt{ДВИГАТЬСЯ: T, (1, 6)} \\
            }
        }
    }
    & & \\ \\

    \theutterance \stepcounter{utterance}  
    & & & \multicolumn{2}{p{0.3\linewidth}}{
        \cellcolor[rgb]{0.9,0.9,0.9}{
            \makecell[{{p{\linewidth}}}]{
                \texttt{\tiny{[GM$|$GM]}}
                \texttt{T moved to (1, 6)} \\
            }
        }
    }
    & & \\ \\

    \theutterance \stepcounter{utterance}  
    & & \multicolumn{4}{p{0.6\linewidth}}{
        \cellcolor[rgb]{0.9,0.9,0.9}{
            \makecell[{{p{\linewidth}}}]{
                \texttt{\tiny{[P1$\langle$GM]}}
                \texttt{Успешно перемещён объект "W" в (7,7). Ваша обновлённая сетка выглядит так:} \\
\texttt{```} \\
\texttt{ 1234567} \\
\texttt{╔═╤═╤═══╗} \\
\texttt{║C│◌│◌◌◌║ 1} \\
\texttt{║◌├─┤◌◌◌║ 2} \\
\texttt{╟─┤◌├─┬─╢ 3} \\
\texttt{║◌│◌│◌│◌║ 4} \\
\texttt{╟─┼─┘◌│◌║ 5} \\
\texttt{║◌│◌H┌┴─╢ 6} \\
\texttt{║◌│I◌│TW║ 7} \\
\texttt{╚═╧══╧══╝} \\
\texttt{```} \\
\texttt{Сейчас вы играете раунд 11 из максимум 20.} \\
\texttt{Вы вместе накопили 18 из 12 штрафов.} \\
\texttt{Другой игрок переместил объект на своей сетке.} \\
\\ 
\\ 
\texttt{Какая твоя следующая команда?} \\
            }
        }
    }
    & \\ \\

    \theutterance \stepcounter{utterance}  
    & \multicolumn{4}{p{0.6\linewidth}}{
        \cellcolor[rgb]{0.9,0.9,0.9}{
            \makecell[{{p{\linewidth}}}]{
                \texttt{\tiny{[P1$\rangle$GM]}}
                \texttt{ДВИГАТЬСЯ: H, (4, 7)} \\
            }
        }
    }
    & & \\ \\

    \theutterance \stepcounter{utterance}  
    & & & \multicolumn{2}{p{0.3\linewidth}}{
        \cellcolor[rgb]{0.9,0.9,0.9}{
            \makecell[{{p{\linewidth}}}]{
                \texttt{\tiny{[GM$|$GM]}}
                \texttt{ДВИГАТЬСЯ: H, (4, 7)} \\
            }
        }
    }
    & & \\ \\

    \theutterance \stepcounter{utterance}  
    & & & \multicolumn{2}{p{0.3\linewidth}}{
        \cellcolor[rgb]{0.9,0.9,0.9}{
            \makecell[{{p{\linewidth}}}]{
                \texttt{\tiny{[GM$|$GM]}}
                \texttt{H moved to (4, 7)} \\
            }
        }
    }
    & & \\ \\

    \theutterance \stepcounter{utterance}  
    & & \multicolumn{4}{p{0.6\linewidth}}{
        \cellcolor[rgb]{0.9,0.9,0.9}{
            \makecell[{{p{\linewidth}}}]{
                \texttt{\tiny{[GM$\rangle$P2]}}
                \texttt{Успешно перемещён объект "T" в (1,6). Ваша обновлённая сетка выглядит так:} \\
\texttt{```} \\
\texttt{ 1234567} \\
\texttt{╔═╤═╤═══╗} \\
\texttt{║◌│◌│◌◌◌║ 1} \\
\texttt{║◌├─┤◌◌◌║ 2} \\
\texttt{╟─┤◌├─┬─╢ 3} \\
\texttt{║◌│◌│◌│C║ 4} \\
\texttt{╟─┼─┘◌│◌║ 5} \\
\texttt{║T│◌◌┌┴─╢ 6} \\
\texttt{║◌│◌I│HW║ 7} \\
\texttt{╚═╧══╧══╝} \\
\texttt{```} \\
\texttt{Сейчас вы играете раунд 12 из максимум 20.} \\
\texttt{Вы вместе накопили 18 из 12 штрафов.} \\
\texttt{Другой игрок переместил объект на своей сетке.} \\
\\ 
\\ 
\texttt{Какая твоя следующая команда?} \\
            }
        }
    }
    & \\ \\

    \theutterance \stepcounter{utterance}  
    & & & \multicolumn{4}{p{0.6\linewidth}}{
        \cellcolor[rgb]{0.9,0.9,0.9}{
            \makecell[{{p{\linewidth}}}]{
                \texttt{\tiny{[GM$\langle$P2]}}
                \texttt{ДВИГАТЬСЯ: I, (4, 7)} \\
            }
        }
    }
     \\ \\

    \theutterance \stepcounter{utterance}  
    & & & \multicolumn{2}{p{0.3\linewidth}}{
        \cellcolor[rgb]{0.9,0.9,0.9}{
            \makecell[{{p{\linewidth}}}]{
                \texttt{\tiny{[GM$|$GM]}}
                \texttt{ДВИГАТЬСЯ: I, (4, 7)} \\
            }
        }
    }
    & & \\ \\

    \theutterance \stepcounter{utterance}  
    & & & \multicolumn{2}{p{0.3\linewidth}}{
        \cellcolor[rgb]{0.9,0.9,0.9}{
            \makecell[{{p{\linewidth}}}]{
                \texttt{\tiny{[GM$|$GM]}}
                \texttt{Штраф: (4,7) не пусто, а содержит "I".} \\
\\ 
\texttt{Вы вместе накопили 19 из 12 штрафов.} \\
\texttt{Пожалуйста, попробуйте ещё раз!} \\
            }
        }
    }
    & & \\ \\

    \theutterance \stepcounter{utterance}  
    & & \multicolumn{4}{p{0.6\linewidth}}{
        \cellcolor[rgb]{0.9,0.9,0.9}{
            \makecell[{{p{\linewidth}}}]{
                \texttt{\tiny{[GM$\rangle$P2]}}
                \texttt{Штраф: (4,7) не пусто, а содержит "I".} \\
\\ 
\texttt{Вы вместе накопили 19 из 12 штрафов.} \\
\texttt{Пожалуйста, попробуйте ещё раз!} \\
            }
        }
    }
    & \\ \\

    \theutterance \stepcounter{utterance}  
    & & & \multicolumn{4}{p{0.6\linewidth}}{
        \cellcolor[rgb]{0.9,0.9,0.9}{
            \makecell[{{p{\linewidth}}}]{
                \texttt{\tiny{[GM$\langle$P2]}}
                \texttt{ДВИГАТЬСЯ: I, (5, 7)} \\
            }
        }
    }
     \\ \\

    \theutterance \stepcounter{utterance}  
    & & & \multicolumn{2}{p{0.3\linewidth}}{
        \cellcolor[rgb]{0.9,0.9,0.9}{
            \makecell[{{p{\linewidth}}}]{
                \texttt{\tiny{[GM$|$GM]}}
                \texttt{ДВИГАТЬСЯ: I, (5, 7)} \\
            }
        }
    }
    & & \\ \\

    \theutterance \stepcounter{utterance}  
    & & & \multicolumn{2}{p{0.3\linewidth}}{
        \cellcolor[rgb]{0.9,0.9,0.9}{
            \makecell[{{p{\linewidth}}}]{
                \texttt{\tiny{[GM$|$GM]}}
                \texttt{Штраф: (5,7) не пусто, а содержит "│".} \\
\\ 
\texttt{Вы вместе накопили 20 из 12 штрафов.} \\
\texttt{Пожалуйста, попробуйте ещё раз!} \\
            }
        }
    }
    & & \\ \\

    \theutterance \stepcounter{utterance}  
    & & \multicolumn{4}{p{0.6\linewidth}}{
        \cellcolor[rgb]{0.9,0.9,0.9}{
            \makecell[{{p{\linewidth}}}]{
                \texttt{\tiny{[GM$\rangle$P2]}}
                \texttt{Штраф: (5,7) не пусто, а содержит "│".} \\
\\ 
\texttt{Вы вместе накопили 20 из 12 штрафов.} \\
\texttt{Пожалуйста, попробуйте ещё раз!} \\
            }
        }
    }
    & \\ \\

    \theutterance \stepcounter{utterance}  
    & & & \multicolumn{4}{p{0.6\linewidth}}{
        \cellcolor[rgb]{0.9,0.9,0.9}{
            \makecell[{{p{\linewidth}}}]{
                \texttt{\tiny{[GM$\langle$P2]}}
                \texttt{ДВИГАТЬСЯ: I, (4, 5)} \\
            }
        }
    }
     \\ \\

    \theutterance \stepcounter{utterance}  
    & & & \multicolumn{2}{p{0.3\linewidth}}{
        \cellcolor[rgb]{0.9,0.9,0.9}{
            \makecell[{{p{\linewidth}}}]{
                \texttt{\tiny{[GM$|$GM]}}
                \texttt{ДВИГАТЬСЯ: I, (4, 5)} \\
            }
        }
    }
    & & \\ \\

    \theutterance \stepcounter{utterance}  
    & & & \multicolumn{2}{p{0.3\linewidth}}{
        \cellcolor[rgb]{0.9,0.9,0.9}{
            \makecell[{{p{\linewidth}}}]{
                \texttt{\tiny{[GM$|$GM]}}
                \texttt{Штраф: (4,5) не пусто, а содержит "┘".} \\
\\ 
\texttt{Вы вместе накопили 21 из 12 штрафов.} \\
\texttt{Пожалуйста, попробуйте ещё раз!} \\
            }
        }
    }
    & & \\ \\

    \theutterance \stepcounter{utterance}  
    & & \multicolumn{4}{p{0.6\linewidth}}{
        \cellcolor[rgb]{0.9,0.9,0.9}{
            \makecell[{{p{\linewidth}}}]{
                \texttt{\tiny{[GM$\rangle$P2]}}
                \texttt{Штраф: (4,5) не пусто, а содержит "┘".} \\
\\ 
\texttt{Вы вместе накопили 21 из 12 штрафов.} \\
\texttt{Пожалуйста, попробуйте ещё раз!} \\
            }
        }
    }
    & \\ \\

    \theutterance \stepcounter{utterance}  
    & & & \multicolumn{4}{p{0.6\linewidth}}{
        \cellcolor[rgb]{0.9,0.9,0.9}{
            \makecell[{{p{\linewidth}}}]{
                \texttt{\tiny{[GM$\langle$P2]}}
                \texttt{ДВИГАТЬСЯ: I, (4, 1)} \\
            }
        }
    }
     \\ \\

    \theutterance \stepcounter{utterance}  
    & & & \multicolumn{2}{p{0.3\linewidth}}{
        \cellcolor[rgb]{0.9,0.9,0.9}{
            \makecell[{{p{\linewidth}}}]{
                \texttt{\tiny{[GM$|$GM]}}
                \texttt{ДВИГАТЬСЯ: I, (4, 1)} \\
            }
        }
    }
    & & \\ \\

    \theutterance \stepcounter{utterance}  
    & & & \multicolumn{2}{p{0.3\linewidth}}{
        \cellcolor[rgb]{0.9,0.9,0.9}{
            \makecell[{{p{\linewidth}}}]{
                \texttt{\tiny{[GM$|$GM]}}
                \texttt{Штраф: (4,1) не пусто, а содержит "│".} \\
\\ 
\texttt{Вы вместе накопили 22 из 12 штрафов.} \\
\texttt{Пожалуйста, попробуйте ещё раз!} \\
            }
        }
    }
    & & \\ \\

    \theutterance \stepcounter{utterance}  
    & & \multicolumn{4}{p{0.6\linewidth}}{
        \cellcolor[rgb]{0.9,0.9,0.9}{
            \makecell[{{p{\linewidth}}}]{
                \texttt{\tiny{[GM$\rangle$P2]}}
                \texttt{Штраф: (4,1) не пусто, а содержит "│".} \\
\\ 
\texttt{Вы вместе накопили 22 из 12 штрафов.} \\
\texttt{Пожалуйста, попробуйте ещё раз!} \\
            }
        }
    }
    & \\ \\

    \theutterance \stepcounter{utterance}  
    & & & \multicolumn{4}{p{0.6\linewidth}}{
        \cellcolor[rgb]{0.9,0.9,0.9}{
            \makecell[{{p{\linewidth}}}]{
                \texttt{\tiny{[GM$\langle$P2]}}
                \texttt{ДВИГАТЬСЯ: I, (2, 1)} \\
            }
        }
    }
     \\ \\

    \theutterance \stepcounter{utterance}  
    & & & \multicolumn{2}{p{0.3\linewidth}}{
        \cellcolor[rgb]{0.9,0.9,0.9}{
            \makecell[{{p{\linewidth}}}]{
                \texttt{\tiny{[GM$|$GM]}}
                \texttt{ДВИГАТЬСЯ: I, (2, 1)} \\
            }
        }
    }
    & & \\ \\

    \theutterance \stepcounter{utterance}  
    & & & \multicolumn{2}{p{0.3\linewidth}}{
        \cellcolor[rgb]{0.9,0.9,0.9}{
            \makecell[{{p{\linewidth}}}]{
                \texttt{\tiny{[GM$|$GM]}}
                \texttt{Штраф: (2,1) не пусто, а содержит "│".} \\
\\ 
\texttt{Вы вместе накопили 23 из 12 штрафов.} \\
\texttt{Пожалуйста, попробуйте ещё раз!} \\
            }
        }
    }
    & & \\ \\

    \theutterance \stepcounter{utterance}  
    & & \multicolumn{4}{p{0.6\linewidth}}{
        \cellcolor[rgb]{0.9,0.9,0.9}{
            \makecell[{{p{\linewidth}}}]{
                \texttt{\tiny{[GM$\rangle$P2]}}
                \texttt{Штраф: (2,1) не пусто, а содержит "│".} \\
\\ 
\texttt{Вы вместе накопили 23 из 12 штрафов.} \\
\texttt{Пожалуйста, попробуйте ещё раз!} \\
            }
        }
    }
    & \\ \\

    \theutterance \stepcounter{utterance}  
    & & & \multicolumn{4}{p{0.6\linewidth}}{
        \cellcolor[rgb]{0.9,0.9,0.9}{
            \makecell[{{p{\linewidth}}}]{
                \texttt{\tiny{[GM$\langle$P2]}}
                \texttt{ДВИГАТЬСЯ: I, (3, 1)} \\
            }
        }
    }
     \\ \\

    \theutterance \stepcounter{utterance}  
    & & & \multicolumn{2}{p{0.3\linewidth}}{
        \cellcolor[rgb]{0.9,0.9,0.9}{
            \makecell[{{p{\linewidth}}}]{
                \texttt{\tiny{[GM$|$GM]}}
                \texttt{ДВИГАТЬСЯ: I, (3, 1)} \\
            }
        }
    }
    & & \\ \\

    \theutterance \stepcounter{utterance}  
    & & & \multicolumn{2}{p{0.3\linewidth}}{
        \cellcolor[rgb]{0.9,0.9,0.9}{
            \makecell[{{p{\linewidth}}}]{
                \texttt{\tiny{[GM$|$GM]}}
                \texttt{I moved to (3, 1)} \\
            }
        }
    }
    & & \\ \\

    \theutterance \stepcounter{utterance}  
    & & \multicolumn{4}{p{0.6\linewidth}}{
        \cellcolor[rgb]{0.9,0.9,0.9}{
            \makecell[{{p{\linewidth}}}]{
                \texttt{\tiny{[P1$\langle$GM]}}
                \texttt{Успешно перемещён объект "H" в (4,7). Ваша обновлённая сетка выглядит так:} \\
\texttt{```} \\
\texttt{ 1234567} \\
\texttt{╔═╤═╤═══╗} \\
\texttt{║C│◌│◌◌◌║ 1} \\
\texttt{║◌├─┤◌◌◌║ 2} \\
\texttt{╟─┤◌├─┬─╢ 3} \\
\texttt{║◌│◌│◌│◌║ 4} \\
\texttt{╟─┼─┘◌│◌║ 5} \\
\texttt{║◌│◌◌┌┴─╢ 6} \\
\texttt{║◌│IH│TW║ 7} \\
\texttt{╚═╧══╧══╝} \\
\texttt{```} \\
\texttt{Сейчас вы играете раунд 12 из максимум 20.} \\
\texttt{Вы вместе накопили 23 из 12 штрафов.} \\
\texttt{Другой игрок переместил объект на своей сетке.} \\
\\ 
\\ 
\texttt{Какая твоя следующая команда?} \\
            }
        }
    }
    & \\ \\

    \theutterance \stepcounter{utterance}  
    & \multicolumn{4}{p{0.6\linewidth}}{
        \cellcolor[rgb]{0.9,0.9,0.9}{
            \makecell[{{p{\linewidth}}}]{
                \texttt{\tiny{[P1$\rangle$GM]}}
                \texttt{ДВИГАТЬСЯ: C, (5, 7)} \\
            }
        }
    }
    & & \\ \\

    \theutterance \stepcounter{utterance}  
    & & & \multicolumn{2}{p{0.3\linewidth}}{
        \cellcolor[rgb]{0.9,0.9,0.9}{
            \makecell[{{p{\linewidth}}}]{
                \texttt{\tiny{[GM$|$GM]}}
                \texttt{ДВИГАТЬСЯ: C, (5, 7)} \\
            }
        }
    }
    & & \\ \\

    \theutterance \stepcounter{utterance}  
    & & & \multicolumn{2}{p{0.3\linewidth}}{
        \cellcolor[rgb]{0.9,0.9,0.9}{
            \makecell[{{p{\linewidth}}}]{
                \texttt{\tiny{[GM$|$GM]}}
                \texttt{Штраф: (5,7) не пусто, а содержит "│".} \\
\\ 
\texttt{Вы вместе накопили 24 из 12 штрафов.} \\
\texttt{Пожалуйста, попробуйте ещё раз!} \\
            }
        }
    }
    & & \\ \\

    \theutterance \stepcounter{utterance}  
    & & \multicolumn{4}{p{0.6\linewidth}}{
        \cellcolor[rgb]{0.9,0.9,0.9}{
            \makecell[{{p{\linewidth}}}]{
                \texttt{\tiny{[P1$\langle$GM]}}
                \texttt{Штраф: (5,7) не пусто, а содержит "│".} \\
\\ 
\texttt{Вы вместе накопили 24 из 12 штрафов.} \\
\texttt{Пожалуйста, попробуйте ещё раз!} \\
            }
        }
    }
    & \\ \\

    \theutterance \stepcounter{utterance}  
    & \multicolumn{4}{p{0.6\linewidth}}{
        \cellcolor[rgb]{0.9,0.9,0.9}{
            \makecell[{{p{\linewidth}}}]{
                \texttt{\tiny{[P1$\rangle$GM]}}
                \texttt{ДВИГАТЬСЯ: C, (5, 4)} \\
            }
        }
    }
    & & \\ \\

    \theutterance \stepcounter{utterance}  
    & & & \multicolumn{2}{p{0.3\linewidth}}{
        \cellcolor[rgb]{0.9,0.9,0.9}{
            \makecell[{{p{\linewidth}}}]{
                \texttt{\tiny{[GM$|$GM]}}
                \texttt{ДВИГАТЬСЯ: C, (5, 4)} \\
            }
        }
    }
    & & \\ \\

    \theutterance \stepcounter{utterance}  
    & & & \multicolumn{2}{p{0.3\linewidth}}{
        \cellcolor[rgb]{0.9,0.9,0.9}{
            \makecell[{{p{\linewidth}}}]{
                \texttt{\tiny{[GM$|$GM]}}
                \texttt{C moved to (5, 4)} \\
            }
        }
    }
    & & \\ \\

    \theutterance \stepcounter{utterance}  
    & & \multicolumn{4}{p{0.6\linewidth}}{
        \cellcolor[rgb]{0.9,0.9,0.9}{
            \makecell[{{p{\linewidth}}}]{
                \texttt{\tiny{[GM$\rangle$P2]}}
                \texttt{Успешно перемещён объект "I" в (3,1). Ваша обновлённая сетка выглядит так:} \\
\texttt{```} \\
\texttt{ 1234567} \\
\texttt{╔═╤═╤═══╗} \\
\texttt{║◌│I│◌◌◌║ 1} \\
\texttt{║◌├─┤◌◌◌║ 2} \\
\texttt{╟─┤◌├─┬─╢ 3} \\
\texttt{║◌│◌│◌│C║ 4} \\
\texttt{╟─┼─┘◌│◌║ 5} \\
\texttt{║T│◌◌┌┴─╢ 6} \\
\texttt{║◌│◌◌│HW║ 7} \\
\texttt{╚═╧══╧══╝} \\
\texttt{```} \\
\texttt{Сейчас вы играете раунд 13 из максимум 20.} \\
\texttt{Вы вместе накопили 24 из 12 штрафов.} \\
\texttt{Другой игрок переместил объект на своей сетке.} \\
\\ 
\\ 
\texttt{Какая твоя следующая команда?} \\
            }
        }
    }
    & \\ \\

    \theutterance \stepcounter{utterance}  
    & & & \multicolumn{4}{p{0.6\linewidth}}{
        \cellcolor[rgb]{0.9,0.9,0.9}{
            \makecell[{{p{\linewidth}}}]{
                \texttt{\tiny{[GM$\langle$P2]}}
                \texttt{ДВИГАТЬСЯ: C, (5, 7)} \\
            }
        }
    }
     \\ \\

    \theutterance \stepcounter{utterance}  
    & & & \multicolumn{2}{p{0.3\linewidth}}{
        \cellcolor[rgb]{0.9,0.9,0.9}{
            \makecell[{{p{\linewidth}}}]{
                \texttt{\tiny{[GM$|$GM]}}
                \texttt{ДВИГАТЬСЯ: C, (5, 7)} \\
            }
        }
    }
    & & \\ \\

    \theutterance \stepcounter{utterance}  
    & & & \multicolumn{2}{p{0.3\linewidth}}{
        \cellcolor[rgb]{0.9,0.9,0.9}{
            \makecell[{{p{\linewidth}}}]{
                \texttt{\tiny{[GM$|$GM]}}
                \texttt{Штраф: (5,7) не пусто, а содержит "│".} \\
\\ 
\texttt{Вы вместе накопили 25 из 12 штрафов.} \\
\texttt{Пожалуйста, попробуйте ещё раз!} \\
            }
        }
    }
    & & \\ \\

    \theutterance \stepcounter{utterance}  
    & & \multicolumn{4}{p{0.6\linewidth}}{
        \cellcolor[rgb]{0.9,0.9,0.9}{
            \makecell[{{p{\linewidth}}}]{
                \texttt{\tiny{[GM$\rangle$P2]}}
                \texttt{Штраф: (5,7) не пусто, а содержит "│".} \\
\\ 
\texttt{Вы вместе накопили 25 из 12 штрафов.} \\
\texttt{Пожалуйста, попробуйте ещё раз!} \\
            }
        }
    }
    & \\ \\

    \theutterance \stepcounter{utterance}  
    & & & \multicolumn{4}{p{0.6\linewidth}}{
        \cellcolor[rgb]{0.9,0.9,0.9}{
            \makecell[{{p{\linewidth}}}]{
                \texttt{\tiny{[GM$\langle$P2]}}
                \texttt{ДВИГАТЬСЯ: C, (6, 7)} \\
            }
        }
    }
     \\ \\

    \theutterance \stepcounter{utterance}  
    & & & \multicolumn{2}{p{0.3\linewidth}}{
        \cellcolor[rgb]{0.9,0.9,0.9}{
            \makecell[{{p{\linewidth}}}]{
                \texttt{\tiny{[GM$|$GM]}}
                \texttt{ДВИГАТЬСЯ: C, (6, 7)} \\
            }
        }
    }
    & & \\ \\

    \theutterance \stepcounter{utterance}  
    & & & \multicolumn{2}{p{0.3\linewidth}}{
        \cellcolor[rgb]{0.9,0.9,0.9}{
            \makecell[{{p{\linewidth}}}]{
                \texttt{\tiny{[GM$|$GM]}}
                \texttt{Штраф: (6,7) не пусто, а содержит "H".} \\
\\ 
\texttt{Вы вместе накопили 26 из 12 штрафов.} \\
\texttt{Пожалуйста, попробуйте ещё раз!} \\
            }
        }
    }
    & & \\ \\

    \theutterance \stepcounter{utterance}  
    & & \multicolumn{4}{p{0.6\linewidth}}{
        \cellcolor[rgb]{0.9,0.9,0.9}{
            \makecell[{{p{\linewidth}}}]{
                \texttt{\tiny{[GM$\rangle$P2]}}
                \texttt{Штраф: (6,7) не пусто, а содержит "H".} \\
\\ 
\texttt{Вы вместе накопили 26 из 12 штрафов.} \\
\texttt{Пожалуйста, попробуйте ещё раз!} \\
            }
        }
    }
    & \\ \\

    \theutterance \stepcounter{utterance}  
    & & & \multicolumn{4}{p{0.6\linewidth}}{
        \cellcolor[rgb]{0.9,0.9,0.9}{
            \makecell[{{p{\linewidth}}}]{
                \texttt{\tiny{[GM$\langle$P2]}}
                \texttt{ДВИГАТЬСЯ: C, (7, 7)} \\
\\ 
\texttt{(Примечание: (7,7) занято 'W', попробую (5, 4) — нет, там '┘'. Попробую (6, 5) — нет, там '◌'?)} \\
\texttt{Посмотрю на сетку: (6,5) — это '◌'.} \\
\\ 
\texttt{ДВИГАТЬСЯ: C, (6, 5)} \\
            }
        }
    }
     \\ \\

    \theutterance \stepcounter{utterance}  
    & & & \multicolumn{2}{p{0.3\linewidth}}{
        \cellcolor[rgb]{0.9,0.9,0.9}{
            \makecell[{{p{\linewidth}}}]{
                \texttt{\tiny{[GM$|$GM]}}
                \texttt{ДВИГАТЬСЯ: C, (7, 7)} \\
\\ 
\texttt{(Примечание: (7,7) занято 'W', попробую (5, 4) — нет, там '┘'. Попробую (6, 5) — нет, там '◌'?)} \\
\texttt{Посмотрю на сетку: (6,5) — это '◌'.} \\
\\ 
\texttt{ДВИГАТЬСЯ: C, (6, 5)} \\
            }
        }
    }
    & & \\ \\

    \theutterance \stepcounter{utterance}  
    & & & \multicolumn{2}{p{0.3\linewidth}}{
        \cellcolor[rgb]{0.9,0.9,0.9}{
            \makecell[{{p{\linewidth}}}]{
                \texttt{\tiny{[GM$|$GM]}}
                \texttt{Invalid format: head is not empty: 'ДВИГАТЬСЯ: C, (7, 7)} \\
\\ 
\texttt{(Примечание: (7,7) занято 'W', попробую (5, 4) — нет, там '┘'. Попробую (6, 5) — нет, там '◌'?)} \\
\texttt{Посмотрю на сетку: (6,5) — это '◌'.} \\
\\ 
\texttt{'} \\
            }
        }
    }
    & & \\ \\

    \theutterance \stepcounter{utterance}  
    & & & \multicolumn{2}{p{0.3\linewidth}}{
        \cellcolor[rgb]{0.9,0.9,0.9}{
            \makecell[{{p{\linewidth}}}]{
                \texttt{\tiny{[GM$|$GM]}}
                \texttt{Maximum number of penalties exceeded} \\
            }
        }
    }
    & & \\ \\

    \theutterance \stepcounter{utterance}  
    & & & \multicolumn{2}{p{0.3\linewidth}}{
        \cellcolor[rgb]{0.9,0.9,0.9}{
            \makecell[{{p{\linewidth}}}]{
                \texttt{\tiny{[GM$|$GM]}}
                \texttt{Initial states:} \\
\texttt{```} \\
\texttt{Player 1:} \\
\texttt{ 1234567} \\
\texttt{╔═╤═╤═══╗} \\
\texttt{║C│◌│◌TI║ 1} \\
\texttt{║◌├─┤W◌◌║ 2} \\
\texttt{╟─┤◌├─┬─╢ 3} \\
\texttt{║◌│◌│◌│◌║ 4} \\
\texttt{╟─┼─┘◌│◌║ 5} \\
\texttt{║◌│◌H┌┴─╢ 6} \\
\texttt{║◌│◌◌│◌◌║ 7} \\
\texttt{╚═╧══╧══╝} \\
\\ 
\\ 
\texttt{Player 2:} \\
\texttt{ 1234567} \\
\texttt{╔═╤═╤═══╗} \\
\texttt{║H│◌│◌◌◌║ 1} \\
\texttt{║◌├─┤◌◌◌║ 2} \\
\texttt{╟─┤◌├─┬─╢ 3} \\
\texttt{║◌│◌│C│◌║ 4} \\
\texttt{╟─┼─┘◌│◌║ 5} \\
\texttt{║◌│◌◌┌┴─╢ 6} \\
\texttt{║T│WI│◌◌║ 7} \\
\texttt{╚═╧══╧══╝} \\
\\ 
\texttt{```} \\
            }
        }
    }
    & & \\ \\

    \theutterance \stepcounter{utterance}  
    & & & \multicolumn{2}{p{0.3\linewidth}}{
        \cellcolor[rgb]{0.9,0.9,0.9}{
            \makecell[{{p{\linewidth}}}]{
                \texttt{\tiny{[GM$|$GM]}}
                \texttt{End states:} \\
\texttt{```} \\
\texttt{Player 1:} \\
\texttt{ 1234567} \\
\texttt{╔═╤═╤═══╗} \\
\texttt{║◌│◌│◌◌◌║ 1} \\
\texttt{║◌├─┤◌◌◌║ 2} \\
\texttt{╟─┤◌├─┬─╢ 3} \\
\texttt{║◌│◌│C│◌║ 4} \\
\texttt{╟─┼─┘◌│◌║ 5} \\
\texttt{║◌│◌◌┌┴─╢ 6} \\
\texttt{║◌│IH│TW║ 7} \\
\texttt{╚═╧══╧══╝} \\
\\ 
\\ 
\texttt{Player 2:} \\
\texttt{ 1234567} \\
\texttt{╔═╤═╤═══╗} \\
\texttt{║◌│I│◌◌◌║ 1} \\
\texttt{║◌├─┤◌◌◌║ 2} \\
\texttt{╟─┤◌├─┬─╢ 3} \\
\texttt{║◌│◌│◌│C║ 4} \\
\texttt{╟─┼─┘◌│◌║ 5} \\
\texttt{║T│◌◌┌┴─╢ 6} \\
\texttt{║◌│◌◌│HW║ 7} \\
\texttt{╚═╧══╧══╝} \\
\\ 
\texttt{```} \\
            }
        }
    }
    & & \\ \\

    \theutterance \stepcounter{utterance}  
    & & & \multicolumn{2}{p{0.3\linewidth}}{
        \cellcolor[rgb]{0.9,0.9,0.9}{
            \makecell[{{p{\linewidth}}}]{
                \texttt{\tiny{[GM$|$GM]}}
                \texttt{* success: False} \\
\texttt{* lose: False} \\
\texttt{* aborted: True} \\
\texttt{{-}{-}{-}{-}{-}{-}{-}} \\
\texttt{* Shifts: 8.00} \\
\texttt{* Max Shifts: 8.00} \\
\texttt{* Min Shifts: 4.00} \\
\texttt{* End Distance Sum: 15.10} \\
\texttt{* Init Distance Sum: 30.73} \\
\texttt{* Expected Distance Sum: 20.95} \\
\texttt{* Penalties: 27.00} \\
\texttt{* Max Penalties: 12.00} \\
\texttt{* Rounds: 12.00} \\
\texttt{* Max Rounds: 20.00} \\
\texttt{* Object Count: 5.00} \\
            }
        }
    }
    & & \\ \\

\end{supertabular}
}

\end{document}
